% !TeX root = ../../Book1.tex
\chapter{Hausdorff 测度、维数与 Frostman 方法}
\label{b1:ch:15}
% 本章摘要（不计入编号）
\begin{chapterabstract}
一条曲线在平面中通常没有面积，却可以具有正的长度；
Cantor 集甚至可能需要非整数的指数才能辨认其大小。
本章以覆盖集直径的幂作为代价，构造 Hausdorff 测度并定义相应的维数。
等直径不等式把整数维的归一化与 Lebesgue 体积联系起来；
质量分布原理和 Frostman 引理则把维数下界改写成测度增长问题。
最后讨论小球密度，为后续 Lipschitz 映射的微分及面积公式准备几何语言。
\end{chapterabstract}

% 本章目标（不计入编号）
\begin{learninggoals}
\begin{itemize}
\item\label{b1:item:ch15-goal-cover}区分固定尺度的覆盖量、无尺度限制的覆盖量与 Hausdorff 测度，掌握外测度和维数的基本性质。
\item\label{b1:item:ch15-goal-normalize}理解等直径不等式及其在整数维测度归一化中的作用，不把体积比较误解为包围球结论。
\item\label{b1:item:ch15-goal-mass}能用覆盖给出维数上界，用质量分布、Frostman 测度或有限能量给出维数下界。
\item\label{b1:item:ch15-goal-density}掌握有限 Hausdorff 测度集的上密度估计，分清上密度控制、密度存在和可参数化结构。
\end{itemize}
\end{learninggoals}

前两章的覆盖定理与测度微分在本章有不同的用途：前者用于比较覆盖代价，
后者提供平面上密度为 \(1\) 的基准。第\ref{b1:ch:12}章的紧性则在 Frostman 测度的构造中重新出现。
首次阅读可以先沿“测度定义—Cantor 维数—质量分布—Frostman 引理”这一顺序掌握维数估计，
再回读等直径不等式及密度证明中的常数。

Fremlin 的《Measure Theory》第264节与 Simon 的
《Lectures on Geometric Measure Theory》前几节%
\cite{Fremlin2016Measure,Simon1983Lectures}分别提供测度构造与几何密度的参考。
本章对所用结果给出证明，尤其写清零测遗漏的处理、紧性构造与端点条件；
需要更多分形实例的读者，可结合 Falconer 的
《Fractal Geometry: Mathematical Foundations and Applications》%
\cite{Falconer2014Fractal}阅读。

% !TeX root = ../../Book1.tex
\section{Hausdorff 测度}
\label{b1:sec:1501}

Lebesgue 测度用体积衡量集合，却不能进一步区分许多零测集：一段曲线、一个曲面和三分 Cantor 集，在足够高维的环境空间中都可能只有零体积。若覆盖一个集合时不再累加覆盖块的体积，而是累加其直径的某个幂，就能让尺度与指数共同参与测量。这个构造只需要距离，因此也适用于没有线性结构的度量空间。

% 实读来源：Fremlin2016Measure，第4卷 471A--E、471J--L，作者正式 TeX 全文。
% 改编：保留一般度量空间；补入 s=0 的计数约定，采用本书的全实指数几何归一化。
% Borel 可测性使用 471C 的分层思想；Borel 包络只证明后文够用的版本，不调用描述集合论。
本节的度量空间构造可与 Fremlin 的《Measure Theory》\cite[471A--E、471J--L]{Fremlin2016Measure}对读。该书使用直径幂的原始归一化；本书稍后乘一个只依赖指数的正常数，以便整数维的测度与通常的长度、面积、体积相衔接。两种约定不会改变零测集和维数，但比较具体测度值时必须保留这个常数。

\subsection{直径覆盖与尺度}
\label{b1:subsec:150101}

固定度量空间 \((X,d)\)。集合 \(U\subseteq X\) 的直径为 \(\operatorname{diam}U=\sup\{d(x,y):x,y\in U\}\)，约定空集直径为零，允许无界集合的直径为无穷。下面的覆盖集不要求开、闭或可测；它们只承担几何覆盖的作用。

\begin{definition}
\label{b1:def:hausdorff-content}
给定 \(s\geq0\)、\(0<\delta\leq\infty\) 及 \(E\subseteq X\)，定义覆盖量
\[
h^s_\delta(E)
=\inf\Bigl\{\sum_{i=1}^{\infty}(\operatorname{diam}U_i)^s:
E\subseteq\bigcup_{i=1}^{\infty}U_i,
\ \operatorname{diam}U_i\leq\delta\Bigr\}.
\]
这里有限覆盖可以补空集写成序列；没有满足条件的可数覆盖时，下确界取 \(+\infty\)。当 \(\delta=\infty\) 时不限制直径。代价的零维约定是：若 \(s=0\)，则每个非空覆盖集的代价为 \(1\)，包括单点；空集的代价始终为 \(0\)。若 \(s>0\)，则使用通常的直径幂。
\symidx{\(h^s_\delta(E)\)：直径受限的原始覆盖量}
\end{definition}

零维约定只规定这个覆盖构造的代价，不是把一般算术中的 \(0^0\) 重新定义。无尺度限制的 \(h^s_\infty\) 在英文文献中称为 Hausdorff content；\foreignidx{Hausdorff content}不同资料也可能把后面引入的归一化常数包含在这一记号中。

尺度越小，可用的覆盖越少，因此
\[
0<\delta_1\leq\delta_2
\quad\Longrightarrow\quad
h^s_{\delta_1}(E)\geq h^s_{\delta_2}(E).
\]
这一步限制很重要。例如，两个不同的点可以一起放进一个无尺度限制的覆盖集，所以 \(h^0_\infty\) 对这个两点集的值仍为 \(1\)。若要求覆盖块的直径小于两点间距，就必须使用两个非空块。只有让尺度趋于零，测量才不能把相距一定距离的部分永久合并。

每个固定的 \(h^s_\delta\) 都是外测度。空集与单调性直接来自定义。为说明可数次可加性，设 \(E=\bigcup_jE_j\)。若 \(\sum_jh^s_\delta(E_j)=\infty\)，不等式自动成立；否则，对每个 \(j\) 取代价不超过 \(h^s_\delta(E_j)+\varepsilon2^{-j}\) 的覆盖，将所有这些覆盖合并，得到
\[
h^s_\delta(E)\leq\sum_jh^s_\delta(E_j)+\varepsilon.
\]
令 \(\varepsilon\downarrow0\) 即可。固定尺度的外测度一般还不是 Borel 测度；上述两点集已经说明，它未必对不交闭集可加。

\subsection{\texorpdfstring{\(\mathcal H^s\)}{Hausdorff 测度}的构造}
\label{b1:subsec:150102}

由尺度的单调性，可以定义极限
\[
h^s(E)=\lim_{\delta\downarrow0}h^s_\delta(E)
=\sup_{\delta>0}h^s_\delta(E).
\]
它仍是外测度。事实上，对任意可数族 \((E_j)\) 及任意 \(\delta>0\)，有
\[
h^s_\delta\left(\bigcup_jE_j\right)
\leq\sum_jh^s_\delta(E_j)
\leq\sum_jh^s(E_j),
\]
再令 \(\delta\downarrow0\)。这里直接核对了可数次可加性，并未使用“一族外测度的上确界总是外测度”这样的断言。

为统一不同指数下的记号，先给出归一化常数。对 \(a>0\)，定义\term[gamma-hanshu]{Gamma 函数}[Gamma function]
\[
\Gamma(a)=\int_0^\infty t^{a-1}e^{-t}\,\mathrm dt.
\]
这个积分严格为正且有限：在 \((0,1)\) 上，\(t^{a-1}\) 可积；在 \([1,\infty)\) 上，指数衰减控制任意固定幂。例如取整数 \(N>a\)，由 \(e^{t/2}\geq(t/2)^N/N!\) 可知被积函数不超过某个常数乘 \(e^{-t/2}\)。本章只用这个积分来定义正常数，不需要 Gamma 函数的进一步理论。特别地，\(\Gamma(1)=1\)。

\begin{definition}
\label{b1:def:hausdorff-measure}
给定 \(s\geq0\)，令
\[
c_s=\frac{\pi^{s/2}}{2^s\Gamma(1+s/2)},
\quad
\mathcal H^s_\delta(E)=c_s h^s_\delta(E),
\quad
\mathcal H^s(E)=c_s h^s(E).
\]
称 \(\mathcal H^s\) 为 \(s\) 维\term[hausdorff-ce-du]{Hausdorff 外测度}[Hausdorff outer measure]；把它限制到 Carathéodory 可测集上得到的测度，称为 \(s\) 维\term[hausdorff-cedu]{Hausdorff 测度}[Hausdorff measure]。下面将证明所有 Borel 集都在其可测域内。讨论 Borel 测度时，我们也把相应的限制记为 \(\mathcal H^s\)。
\symidx{\(\mathcal H^s\)：归一化的 Hausdorff 测度及其外测度}
\symidx{\(\mathcal H^s_\delta\)：归一化的尺度覆盖量}
\symidx{\(c_s\)：Hausdorff 测度的归一化常数}
\end{definition}

对任意集合写 \(\mathcal H^s(E)\) 时，尚未声明其可测性也有意义，此时指外测度值。整数 \(d\geq1\) 时，常数将满足 \(c_d=\omega_d/2^d\)，其中 \(\omega_d\) 是 \(\mathbb R^d\) 中单位球的体积；这一识别及 \(\mathcal H^d=m_d\) 的证明放在第\ref{b1:sec:1503}节。当前构造只需要 \(0<c_s<\infty\) 与 \(c_0=1\)。

\begin{theorem}
\label{b1:thm:hausdorff-borel-measure}
任意度量空间上的 \(\mathcal H^s\) 都使全部 Borel 集 Carathéodory 可测。因而它在 Borel \(\sigma\)-代数上的限制是测度，在全部 Carathéodory 可测集上的限制则是完备测度。
\end{theorem}
\begin{proof}
先证明一个由距离带来的可加性。若 \(A,B\subseteq X\) 且 \(\operatorname{dist}(A,B)>0\)，取 \(\delta<\operatorname{dist}(A,B)\)。任何直径不超过 \(\delta\) 的覆盖块都不能同时接触 \(A\) 与 \(B\)。把覆盖 \(A\cup B\) 的块按它所接触的集合分开，得到
\[
h^s_\delta(A\cup B)\geq h^s_\delta(A)+h^s_\delta(B).
\]
令 \(\delta\downarrow0\)，再结合外测度的次可加性，得
\[
\mathcal H^s(A\cup B)=\mathcal H^s(A)+\mathcal H^s(B).
\]
若一个度量空间上的外测度对所有正距离相隔的两集合都满足这个等式，则称其为\term[duliang-waicedu]{度量外测度}[metric outer measure]。

下面证明任意度量外测度 \(\theta\) 都使闭集可测。固定闭集 \(F\)，只需考察满足 \(\theta(E)<\infty\) 的任意集合 \(E\)：若 \(\theta(E)=\infty\)，Carathéodory 判据所需的反向不等式自动成立。空的或全空间的 \(F\) 无需证明，以下设 \(F\) 非空。记
\[
t(x)=\operatorname{dist}(x,F),
\quad E_n=\{x\in E:t(x)\geq1/n\}.
\]
集合 \(E_n\) 与 \(E\cap F\) 正距离相隔，所以
\[
\theta(E)\geq\theta(E_n)+\theta(E\cap F).
\]
不能在尚未证明可测性时直接对外测度使用从下连续性。为补上这一步，令
\[
A_i=\{x\in E:1/(i+1)<t(x)\leq1/i\},
\quad i\geq1.
\]
函数 \(t\) 是 \(1\)-Lipschitz 的。两个下标至少相差 \(2\) 的带 \(A_i\) 正距离相隔，因此任意有限个奇数下标带的测度之和不超过 \(\theta(E)\)，偶数下标也一样。于是
\[
\sum_{i=1}^{\infty}\theta(A_i)\leq2\theta(E)<\infty.
\]
又因 \(F\) 闭，\(E\setminus F\subseteq E_n\cup\bigcup_{i\geq n}A_i\)，故
\[
\theta(E\setminus F)
\leq\theta(E_n)+\sum_{i\geq n}\theta(A_i).
\]
尾和趋于零，结合前面的分离估计，得到
\[
\theta(E\cap F)+\theta(E\setminus F)\leq\theta(E).
\]
闭集因此 Carathéodory 可测。由定理\ref{b1:thm:caratheodory}，可测集组成 \(\sigma\)-代数，包含闭集便包含全部 Borel 集；同一定理也给出可数可加性与完备性。取 \(\theta=\mathcal H^s\) 即得结论。
\end{proof}

\subsection{基本性质}
\label{b1:subsec:150103}

先看指数为零的情形。它检验了单点代价的约定，也说明 Hausdorff 测度并非总是局部有限。

\begin{proposition}
\label{b1:prop:hausdorff-counting}
对任意 \(E\subseteq X\)，\(\mathcal H^0(E)\) 等于 \(E\) 的点数；若 \(E\) 无限，则其值为 \(+\infty\)。若 \(s>0\)，则每个可数集的 \(\mathcal H^s\) 外测度为零。
\end{proposition}
\begin{proof}
若 \(E\) 恰有 \(N\) 个点，单点覆盖给出 \(h^0_\delta(E)\leq N\)。当 \(N\geq2\) 时，取 \(\delta\) 小于这些点之间的最小距离，每个覆盖块至多包含一个点，因此 \(h^0_\delta(E)\geq N\)；空集与单点情形直接成立。由于 \(c_0=1\)，得到 \(\mathcal H^0(E)=N\)。无限集包含任意大的有限子集，单调性遂给出无穷值。若 \(s>0\)，每个单点的代价为零，可数个单点本身便是任意尺度下的零代价覆盖。
\end{proof}

例如，\(\mathbb R\) 上的 \(\mathcal H^0\) 在任何非空开区间上都为无穷。故 Borel 可测性不意味着局部有限，更不能据此直接把任意 Hausdorff 测度称为第\ref{b1:ch:11}章意义下的 Radon 测度。另一方面，覆盖允许使用任意集合并不会妨碍从外部找到一个具有相同测度的 Borel 集。

\begin{proposition}
\label{b1:prop:hausdorff-borel-hull}
对任意 \(E\subseteq X\) 及 \(s\geq0\)，存在 Borel 集 \(G\supseteq E\)，使 \(\mathcal H^s(G)=\mathcal H^s(E)\)。此外，
\[
\mathcal H^s_\infty(E)=0
\quad\Longleftrightarrow\quad
\mathcal H^s(E)=0.
\]
\end{proposition}
\begin{proof}
若 \(h^s(E)=\infty\)，取 \(G=X\)。否则，对每个 \(n\geq1\)，取覆盖 \((U_{ni})_i\)，使各块直径不超过 \(1/n\)，总代价不超过 \(h^s_{1/n}(E)+1/n\)。集合的闭包不改变直径，也不改变是否为空，故
\[
G=\bigcap_{n=1}^{\infty}\bigcup_{i=1}^{\infty}\overline{U_{ni}}
\]
是包含 \(E\) 的 Borel 集。固定 \(\delta>0\)，对一切充分大的 \(n\)，第 \(n\) 组闭包是 \(G\) 的合法 \(\delta\)-覆盖。因此
\[
h^s_\delta(G)\leq h^s_{1/n}(E)+1/n\leq h^s(E)+1/n.
\]
先令 \(n\to\infty\)，再令 \(\delta\downarrow0\)，并用单调性，即得等测度包络。

零测性等价的反向蕴含由 \(h^s_\infty\leq h^s\) 得到。正向蕴含在 \(s=0\) 时只涉及空集。若 \(s>0\) 且 \(h^s_\infty(E)=0\)，给定 \(\varepsilon,\delta>0\)，可以找到总代价小于 \(\min\{\varepsilon,\delta^s\}\) 的无尺度限制覆盖。每个非空块的直径都小于 \(\delta\)，故它也是一个 \(\delta\)-覆盖，且 \(h^s_\delta(E)<\varepsilon\)。令 \(\varepsilon\downarrow0\)，再对尺度取上确界即可。
\end{proof}

最后一个等价关系并不是说无尺度限制的覆盖量与测度总相等。它只说明，若总代价已经能够任意小，那么覆盖块也被迫同时变小。这正是后面用一次覆盖估计判断零测性时所需的事实。

\subsection{Lipschitz 映射下的行为}
\label{b1:subsec:150104}

Hausdorff 测度最直接的几何性质来自直径的控制。若映射使所有距离至多放大 \(L\) 倍，就应当使 \(s\) 维代价至多放大 \(L^s\) 倍。这里不需要映射可微，也不要求集合先有某种曲面参数化。

\begin{proposition}
\label{b1:prop:hausdorff-lipschitz}
设 \(X,Y\) 为度量空间，\(A\subseteq X\)，映射 \(f:A\to Y\) 满足
\[
d_Y\bigl(f(x),f(y)\bigr)\leq Ld_X(x,y)
\quad(x,y\in A),
\]
其中 \(L>0\)。则对任意 \(E\subseteq A\) 及 \(s\geq0\)，
\[
\mathcal H^s_Y\bigl(f(E)\bigr)\leq L^s\mathcal H^s_X(E).
\]
Hausdorff 外测度不依赖于所采用的等距环境空间。特别地，等距映射保持 \(\mathcal H^s\)；若 \(f\) 把所有距离恰好乘以 \(L>0\)，则上式为等式。
\end{proposition}
\begin{proof}
先说明子空间问题。若 \(A\subseteq X\)，则 \(A\) 内的覆盖也是 \(X\) 内的覆盖；反过来，把 \(X\) 内的每个覆盖块 \(U_i\) 换成 \(U_i\cap A\)，直径和代价均不增。这证明对 \(E\subseteq A\)，每个尺度的覆盖量在 \(A\) 与 \(X\) 中相同。

现在取 \(E\) 的一个 \(\delta\)-覆盖 \((U_i)\)，把覆盖块先与 \(E\) 相交再取像。它们覆盖 \(f(E)\)，且
\[
\operatorname{diam}\bigl(f(U_i\cap E)\bigr)
\leq L\operatorname{diam}U_i.
\]
当 \(s>0\) 时直接对直径取幂；当 \(s=0\) 时，非空像的数量不超过非空原覆盖块的数量。因此
\[
h^s_{L\delta,Y}\bigl(f(E)\bigr)\leq L^s h^s_{\delta,X}(E).
\]
令 \(\delta\downarrow0\)，乘以相同的 \(c_s\)，即得测度不等式。等距情形分别应用于 \(f\) 与其在像上的逆映射；距离恰好乘以 \(L\) 的情形，再对逆映射使用常数 \(1/L\)，便得到反向不等式。
\end{proof}

若 Lipschitz 常数为零，映射的像至多只有一个点：对 \(s>0\)，像的测度为零；对 \(s=0\)，仍有像的点数不超过原集合点数。这样单独处理，就不必在测度无穷时解释 \(0\cdot\infty\)。欧氏空间中的平移、旋转保持距离，伸缩把距离乘以伸缩率，因此上述结论已经包含了 Hausdorff 测度相应的不变性与齐次性。

% !TeX root = ../../Book1.tex
\section{Hausdorff 维数}
\label{b1:sec:1502}

上一节对每个指数 \(s\) 都构造了一个测度。现在固定集合而改变指数：同一批小覆盖块，在较大的指数下有较小的代价。这个比较会产生一个临界值。在临界值以下，覆盖代价积累为无穷；在临界值以上，覆盖代价可以趋于零。Hausdorff 维数就是从这族测度中提取出的临界指数。

% 实读来源：Fremlin2016Measure，第4卷 471L、471Xh，第2卷 264J，作者正式 TeX 全文。
% 471Xh 只给维数与可数稳定性的习题题面，以下补写完整证明。
% Cantor 部分推广为等比例两分支；下界另用二进位推前测度和球质量估计，非逐句改写 264J 的精确测度证明。

\subsection{临界指数}
\label{b1:subsec:150201}

\begin{proposition}
\label{b1:prop:hausdorff-critical-exponent}
设 \(E\) 是度量空间的任意子集，\(0\leq s<t\)。对每个 \(0<\delta<\infty\)，有
\[
h^t_\delta(E)\leq\delta^{t-s}h^s_\delta(E),
\quad
\mathcal H^t_\delta(E)
\leq\frac{c_t}{c_s}\delta^{t-s}\mathcal H^s_\delta(E).
\]
因此，若 \(\mathcal H^s(E)<\infty\)，则 \(\mathcal H^t(E)=0\)；等价地，若 \(\mathcal H^t(E)>0\)，则 \(\mathcal H^s(E)=\infty\)。
\end{proposition}
\begin{proof}
对任意合法的 \(\delta\)-覆盖，当 \(s>0\) 时逐项有
\[
(\operatorname{diam}U_i)^t
\leq\delta^{t-s}(\operatorname{diam}U_i)^s.
\]
当 \(s=0\) 时，每个非空块的右端是 \(\delta^t\)，不等式同样成立；空块两端均为零。求和并取下确界，得到尺度估计。若 \(h^s(E)<\infty\)，则
\[
h^t_\delta(E)\leq\delta^{t-s}h^s(E)\longrightarrow0
\quad(\delta\downarrow0),
\]
故 \(\mathcal H^t(E)=0\)。最后的结论是这一蕴含的逆否命题。
\end{proof}

必须让尺度趋零，才会得到这种明确的临界行为。对于一个任意大的覆盖块，增大指数未必减小代价；几何归一化常数 \(c_s\) 也随 \(s\) 改变。因此本命题不是在说 \(\mathcal H^s(E)\) 对指数处处单调，而是在比较不同指数下有限、无穷与零这三种状态。特别是，一旦某个指数给出有限测度，所有更大的指数便只得到零测度。

\subsection{Hausdorff 维数}
\label{b1:subsec:150202}

\begin{definition}
\label{b1:def:hausdorff-dimension}
给定度量空间中的集合 \(E\)，定义其\term[hausdorff-weishu]{Hausdorff 维数}[Hausdorff dimension]为
\[
\dim_{\mathrm H}E
=\inf\{s\geq0:\mathcal H^s(E)=0\},
\]
其中空集的下确界取 \(+\infty\)。这里 \(\dim_{\mathrm H}\varnothing=0\)。
\symidx{\(\dim_{\mathrm H}E\)：集合的 Hausdorff 维数}
\end{definition}

若 \(d_E=\dim_{\mathrm H}E\)，则
\[
\mathcal H^s(E)=\infty\quad(0\leq s<d_E),
\quad
\mathcal H^t(E)=0\quad(t>d_E).
\]
第二个断言在 \(d_E<\infty\) 时由下确界定义证明：给定 \(t>d_E\)，可取 \(u<t\) 使 \(\mathcal H^u(E)=0\)，再用命题\ref{b1:prop:hausdorff-critical-exponent}。若第一个断言失败，即某个 \(s<d_E\) 的测度有限，则在 \(s<t<d_E\) 处测度已经为零，与 \(d_E\) 的定义矛盾。这也说明
\[
\dim_{\mathrm H}E
=\sup\{s\geq0:\mathcal H^s(E)=\infty\},
\]
其中右侧集合为空时取值为零。定义没有决定临界指数处的测度值；知道维数之后，这个值仍需要另外计算。

例如，所有非空有限集和所有可数无限集的维数都为零，但它们的 \(\mathcal H^0\) 分别是有限点数与无穷。维数相同不等于测度相同。另一方面，由单调性立即得到
\[
E\subseteq F\quad\Longrightarrow\quad
\dim_{\mathrm H}E\leq\dim_{\mathrm H}F.
\]

Lipschitz 映射不增大 Hausdorff 维数。事实上，若 \(t>\dim_{\mathrm H}E\)，则 \(\mathcal H^t(E)=0\)，命题\ref{b1:prop:hausdorff-lipschitz}给出 \(\mathcal H^t\bigl(f(E)\bigr)=0\)，再令 \(t\) 下降到原集合的维数即可；常值映射单独由像至多为单点处理。给定单射 \(f:E\to Y\)，若存在 \(0<a\leq b<\infty\) 满足
\[
a\,d_X(x,y)\leq d_Y\bigl(f(x),f(y)\bigr)
\leq b\,d_X(x,y)
\quad(x,y\in E),
\]
则称 \(f\) 为\term[shuang-lipschitz-yingshe]{双 Lipschitz 映射}[bi-Lipschitz map]。它与在像上的逆映射都是 Lipschitz 的，因而保持 Hausdorff 维数。这种不变性依赖距离的定量比较，仅仅知道一个映射及其逆连续，还不足以使用上述估计。

\subsection{可数稳定性}
\label{b1:subsec:150203}

\begin{proposition}
\label{b1:prop:hausdorff-dimension-countable}
设 \((E_j)_{j\geq1}\) 是同一度量空间中的任意集合序列，则
\[
\dim_{\mathrm H}\left(\bigcup_{j=1}^{\infty}E_j\right)
=\sup_{j\geq1}\dim_{\mathrm H}E_j.
\]
\end{proposition}
\begin{proof}
记右端为 \(a\)。每个 \(E_j\) 都包含在并集中，故左端至少为 \(a\)。若 \(a=\infty\)，证明已经结束；否则，对任意 \(t>a\)，每个 \(E_j\) 都满足 \(\mathcal H^t(E_j)=0\)。外测度的可数次可加性给出
\[
\mathcal H^t\left(\bigcup_jE_j\right)
\leq\sum_j\mathcal H^t(E_j)=0.
\]
所以并集的维数不超过任意 \(t>a\)，从而不超过 \(a\)。
\end{proof}

这个结果允许先处理有界部分，再恢复整个欧氏空间。它也解释了为什么可数集的维数都是零：可数并不会累积单点的维数。若把可数并换成任意并，结论便不成立，因为每个集合都是其单点子集之并。

\begin{proposition}
\label{b1:prop:euclidean-hausdorff-dimension}
对任意 \(E\subseteq\mathbb R^d\)，有 \(\dim_{\mathrm H}E\leq d\)。若其 Lebesgue 外测度 \(m_d^*(E)>0\)，则 \(\dim_{\mathrm H}E=d\)。特别地，非空开集的 Hausdorff 维数为 \(d\)。
\end{proposition}
\begin{proof}
先设 \(E\) 有界，把它放在 \([-M,M]^d\) 内。边长为 \(1/n\) 的网格立方体中，至多 \(C_{M,d}n^d\) 个就能覆盖这个盒子，而每个立方体的直径为 \(\sqrt d/n\)。若 \(t>d\)，这些覆盖的总代价不超过
\[
C_{M,d}n^d\left(\frac{\sqrt d}{n}\right)^t
=C_{M,d}d^{t/2}n^{d-t}\longrightarrow0.
\]
因此 \(\mathcal H^t(E)=0\)。一般集合写成 \(E=\bigcup_{M=1}^{\infty}\bigl(E\cap[-M,M]^d\bigr)\)，由可数稳定性得到维数上界。

下界不必预先使用下一节的精确归一化。若 \(U\) 的有限直径为 \(D>0\)，选 \(x\in U\)，则 \(U\) 包含在以 \(x\) 为中心、边长 \(2D\) 的立方体中，故 \(m_d^*(U)\leq2^dD^d\)；若 \(D=0\)，集合至多为单点，仍满足这个估计。对任意有限尺度的覆盖求和，得到
\[
m_d^*(E)\leq2^d h^d_\delta(E),
\quad
\mathcal H^d(E)\geq c_d2^{-d}m_d^*(E).
\]
正 Lebesgue 外测度因此推出 \(\mathcal H^d(E)>0\)，再由临界指数性质得到 \(\dim_{\mathrm H}E\geq d\)。非空开集包含一个非退化立方体，所以具有正 Lebesgue 测度。
\end{proof}

由此可见，取闭包可能改变维数：\(\mathbb Q^d\) 的维数为零，闭包 \(\mathbb R^d\) 的维数却为 \(d\)。覆盖的尺度虽然趋于零，却不意味着维数只由集合的闭包决定。相反，零体积也还不能决定维数是否小于环境维数；下面先给出严格位于零与一之间的具体值。

\subsection{Cantor 型集合}
\label{b1:subsec:150204}

命题\ref{b1:prop:cantor-null}已经证明三分 Cantor 集不可数而 Lebesgue 测度为零。现在考虑一族类似集合。固定 \(0<\lambda<1/2\)，从 \([0,1]\) 开始，在每个保留区间的左右两端各保留原长度 \(\lambda\) 倍的闭区间。记第 \(n\) 轮的并集为 \(K_{\lambda,n}\)，并令
\[
K_\lambda=\bigcap_{n=0}^{\infty}K_{\lambda,n}.
\]
第 \(n\) 轮恰有 \(2^n\) 个长度为 \(\lambda^n\) 的互不相交闭区间，以下称它们为第 \(n\) 层区间。集合 \(K_\lambda\) 是非空紧集，且
\[
m_1(K_\lambda)\leq m_1(K_{\lambda,n})=(2\lambda)^n\longrightarrow0.
\]

直接用第 \(n\) 层覆盖很容易得到维数上界。困难是下界：定义允许任意覆盖，不能只证明这一组自然覆盖的代价较大。为控制所有覆盖，我们在集合上放置一个概率测度，并证明小球容纳的质量不能太多。

\begin{lemma}[质量分布原理]
\label{b1:lem:mass-distribution}
设 \(\nu\) 为度量空间 \(X\) 上的有限正 Borel 测度，\(E\subseteq X\) 为 Borel 集且 \(\nu(E)>0\)。若存在 \(s>0\)、\(C>0\)、\(r_0>0\)，使
\[
\nu\bigl(\overline B(x,r)\bigr)\leq Cr^s
\quad(x\in E,\ 0<r\leq r_0),
\]
则
\[
\mathcal H^s(E)\geq\frac{c_s\nu(E)}{C}>0,
\quad
\dim_{\mathrm H}E\geq s.
\]
这一估计称为\term[zhiliang-fenbu-yuanli]{质量分布原理}[mass distribution principle]。
\end{lemma}
\begin{proof}
对 \(x\in E\)，令 \(r\downarrow0\)，由有限测度的从上连续性得 \(\nu(\{x\})=0\)。固定 \(0<\delta\leq r_0\)，取 \(E\) 的任意 \(\delta\)-覆盖 \((U_i)\)，丢掉不与 \(E\) 相交的块，并在其余每块中选 \(x_i\in U_i\cap E\)。若 \(D_i=\operatorname{diam}U_i>0\)，则
\[
U_i\cap E\subseteq\overline B(x_i,D_i),
\quad
\nu\bigl(\overline B(x_i,D_i)\bigr)\leq CD_i^s.
\]
若 \(D_i=0\)，这一块至多为一个 \(\nu\)-零测单点。这些闭球及单点覆盖 \(E\)，故由可数次可加性，
\[
\nu(E)\leq C\sum_i(\operatorname{diam}U_i)^s.
\]
取覆盖的下确界，再令尺度趋零并乘以 \(c_s\)，即得结论。证明没有要求原覆盖块可测，也不要求球估计在大尺度上成立。
\end{proof}

\begin{proposition}
\label{b1:prop:cantor-hausdorff-dimension}
对 \(0<\lambda<1/2\)，令 \(s_\lambda=\log2/\log(1/\lambda)\)。则
\[
0<\mathcal H^{s_\lambda}(K_\lambda)<\infty,
\quad
\dim_{\mathrm H}K_\lambda=s_\lambda.
\]
特别地，三分 Cantor 集的 Hausdorff 维数为 \(\log2/\log3\)。
\end{proposition}
\begin{proof}
简记 \(s=s_\lambda\)，于是 \(\lambda^s=1/2\)。第 \(n\) 层区间给出的覆盖满足
\[
h^s_{\lambda^n}(K_\lambda)
\leq2^n(\lambda^n)^s=1.
\]
令 \(n\to\infty\)，得到 \(\mathcal H^s(K_\lambda)\leq c_s\)。

现在构造均匀分配在两分支上的概率测度。对 \(t\in[0,1)\)，定义二进位数字
\[
\varepsilon_j(t)=\lfloor2^jt\rfloor-2\lfloor2^{j-1}t\rfloor\in\{0,1\},
\quad
\pi_\lambda(t)=(1-\lambda)\sum_{j=1}^{\infty}\varepsilon_j(t)\lambda^{j-1}.
\]
这个级数收敛，且其部分和可测，所以 \(\pi_\lambda\) 是 Borel 可测映射。前 \(n\) 位确定第 \(n\) 层的一个区间，后面的尾和位于长度 \(\lambda^n\) 的该区间内，故 \(\pi_\lambda(t)\in K_\lambda\)。将 \([0,1)\) 上的 Lebesgue 测度推前，得到
\[
\nu_\lambda(A)=m_1\bigl(\pi_\lambda^{-1}(A)\bigr)
\quad(A\in\mathcal B(\mathbb R)).
\]
它是集中在 \(K_\lambda\) 上的 Borel 概率测度。每组前 \(n\) 位对应一个长度为 \(2^{-n}\) 的半开二进区间；不同组对应的第 \(n\) 层区间正距离相隔，因此每个第 \(n\) 层区间的 \(\nu_\lambda\) 测度恰为 \(2^{-n}\)。二进端点已经由上述取整公式固定，不影响质量计算。

取 \(0<r<1\)，选 \(n\geq1\) 使 \(\lambda^n\leq r<\lambda^{n-1}\)。任意两个相邻的第 \(n\) 层区间之间，距离至少为
\[
g_n=(1-\lambda)\lambda^{n-1}-\lambda^n
=(1-2\lambda)\lambda^{n-1}.
\]
若一个长度为 \(2r\) 的闭球与 \(N\) 个第 \(n\) 层区间相交，在各交集中选一点并按大小排列，相邻所选点之间至少相隔 \(g_n\)。所以
\[
(N-1)g_n\leq2r<2\lambda^{n-1},
\quad
N\leq1+\frac{2}{1-2\lambda}=:C_\lambda.
\]
由各层区间的质量公式，对任意中心 \(x\in\mathbb R\)，
\[
\nu_\lambda\bigl(\overline B(x,r)\bigr)
\leq C_\lambda2^{-n}
=C_\lambda(\lambda^n)^s
\leq C_\lambda r^s.
\]
当 \(r=1\) 时也有这个估计，因为总质量为 \(1\)。质量分布原理因而给出
\[
\frac{c_s}{C_\lambda}\leq\mathcal H^s(K_\lambda)\leq c_s.
\]
临界指数性质随即推出 \(\dim_{\mathrm H}K_\lambda=s\)。
\end{proof}

这里维数的数值来自两种数量的平衡：每细分一层，覆盖块的数目乘以 \(2\)，直径乘以 \(\lambda\)，所以临界指数满足 \(2\lambda^s=1\)。上界使用这种平衡，下界则还使用分支之间的间隙；没有后者，不能仅凭“块数乘直径幂”的计算断定维数。

若希望继续计算三分 Cantor 集在原始归一化下的精确测度值，可阅读 Fremlin 的《Measure Theory》\cite[264J]{Fremlin2016Measure}。本节保留足以确定维数的上下界，使同一证明能处理全部 \(0<\lambda<1/2\)。下一节转向欧氏空间的精确测度识别；第\ref{b1:sec:1504}节则进一步讨论，什么时候可以从集合的覆盖性质反过来构造具有球质量控制的测度。

% !TeX root = ../../Book1.tex
\section{\texorpdfstring{\(\mathcal H^d\)}{Hd} 与 Lebesgue 测度}
\label{b1:sec:1503}

% 实际读取：Fremlin2016Measure，作者官方 chap26.tex，264H--I 完整证明及264节末注释。
% 264H的有限次坐标对称化在此完整展开；不调用Brunn--Minkowski或面积公式。
% 264I的球覆盖比较改接本书13.1 Vitali，另由第三章立方体覆盖补全零测遗漏。
% 归一化常数以本书070504已证一维Gauss积分、Tonelli和Gamma定义重新推导。

用直径计算覆盖的代价，是否仍能得到通常的长度、面积和体积？
仅凭平移不变性还不能确定答案中的常数。
关键是一个几何问题：给定直径时，集合的体积至多有多大。
本节先证明球达到最大的体积，再把这一不等式与 Vitali 覆盖结合，
精确比较 \(\mathbb R^d\) 上的 \(d\) 维 Hausdorff 测度与 Lebesgue 测度。

\subsection{等直径问题}
\label{b1:subsec:150301}

沿用第\ref{b1:sec:1303}节的记号
\(\omega_d=m_d\bigl(B(0,1)\bigr)\)，其中 \(d\geq1\)。
球面的 Lebesgue 测度为零，所以使用开单位球或闭单位球得到同一个 \(\omega_d\)。
若一个集合的直径为 \(D\)，任选其中一点为中心，只能直接把它放进半径 \(D\) 的球。
要把体积估计改善到半径 \(D/2\) 的球，不能单靠这种包含关系。

\begin{theorem}[等直径不等式]\label{b1:thm:isodiametric}
设 \(A\subseteq\mathbb R^d\) 有界。则
\begin{equation}\label{b1:eq:isodiametric}
 m_d^*(A)\leq\frac{\omega_d}{2^d}
                  (\operatorname{diam}A)^d.
\end{equation}
空集的直径按零理解。对任意 \(D>0\)，直径为 \(D\) 的闭球达到等号。
\end{theorem}

这就是\term[dengzhijing-budengshi]{等直径不等式}[isodiametric inequality]。
它不要求 \(A\) 可测，因此以后可以用于 Hausdorff 覆盖中任意形状的覆盖集。
不等式比较的是体积，而不是说每个直径为 \(D\) 的集合都能放进半径 \(D/2\) 的球。
例如边长为 \(1\) 的正三角形的三个顶点，直径为 \(1\)，却不能放进半径 \(1/2\) 的圆盘：
若一个这样的圆盘含两个相距 \(1\) 的顶点，它的中心必须是二者中点，
第三个顶点距该中点为 \(\sqrt3/2>1/2\)。

下面按 Fremlin 的《Measure Theory》\cite[264H及第264节注释]{Fremlin2016Measure}，
把集合逐坐标对称化。每一步都保持或增大体积，同时不增加直径。
只需 \(d\) 步，就能得到一个关于原点中心对称的集合。

\subsection{欧氏球的极值性质}
\label{b1:subsec:150302}

对第 \(j\) 个坐标，记反射
\[
 R_j(x_1,\ldots,x_d)
 =(x_1,\ldots,x_{j-1},-x_j,x_{j+1},\ldots,x_d).
\]
以下操作把平行于某条坐标轴的截面移到该轴坐标为零的位置，
并用同长度的中心区间替代它。
这种操作称为\term[steiner-duichenhua]{Steiner 对称化}[Steiner symmetrization]；
这里仅需它沿坐标方向的一个版本。

\begin{lemma}[逐坐标对称化]\label{b1:lem:coordinate-symmetrization}
设 \(d\geq2\)、\(C\subseteq\mathbb R^d\) 有界，并且对每个 \(i\in J\subseteq\{1,\ldots,d\}\) 都有 \(R_i(C)=C\)。
若 \(j\notin J\)，则存在有界 Borel 集 \(D\)，使
\[
 m_d(D)\geq m_d^*(C),\quad
 \operatorname{diam}D\leq\operatorname{diam}C,
 \quad R_i(D)=D\quad(i\in J\cup\{j\}).
\]
\end{lemma}
\begin{proof}
置换坐标不会改变距离或 Lebesgue 测度，故只需证明 \(j=d\) 的情形。
空集取 \(D=\varnothing\)。以下令 \(F=\overline C\)，它是紧集，满足
\[
 \operatorname{diam}F=\operatorname{diam}C,
 \quad m_d(F)\geq m_d^*(C).
\]
反射与取闭包可交换，所以 \(F\) 仍保留所有 \(J\) 中的坐标对称性。
对 \(y\in\mathbb R^{d-1}\)，令
\[
 F_y=\{\,t\in\mathbb R:(y,t)\in F\,\},\quad
 \ell(y)=m_1(F_y),\quad
 D=\{\,(y,t):|t|<\ell(y)/2\,\}.
\]
由截面测度的可测性，\(\ell\) 为 Borel 函数，故 \(D\) 是 Borel 集。
紧集 \(F\) 的投影有界，各个截面的长度也有统一上界，因此 \(D\) 有界。
每个 \(D_y\) 的长度恰为 \(\ell(y)\)，包括 \(\ell(y)=0\) 时的空截面。
Tonelli 定理给出
\[
 m_d(D)=\int_{\mathbb R^{d-1}}\ell(y)\,\mathrm dy=m_d(F).
\]

接着验证直径。任取 \((y,t),(y',t')\in D\)，这两个对应的 \(F\)-截面都非空。
记它们的最小、最大点分别为
\[
 a=\min F_y,\quad b=\max F_y,
 \quad a'=\min F_{y'},\quad b'=\max F_{y'}.
\]
这些端点属于截面，因为截面紧。又有 \(\ell(y)\leq b-a\)、\(\ell(y')\leq b'-a'\)，所以
\[
 \begin{aligned}
 |t-t'|
 &\leq|t|+|t'|
 <\frac{\ell(y)+\ell(y')}{2}\\
 &\leq\frac{(b-a)+(b'-a')}{2}
 =\frac{(b-a')+(b'-a)}2\\
 &\leq\max\{\,|b-a'|,|b'-a|\,\}.
 \end{aligned}
\]
因此可以从 \((b,a')\) 与 \((a,b')\) 中选一对端点 \((u,v)\)，使 \(|u-v|\geq|t-t'|\)。
两点 \((y,u),(y',v)\) 都属于 \(F\)，从而
\[
 |(y,t)-(y',t')|^2
 =|y-y'|^2+|t-t'|^2
 \leq|y-y'|^2+|u-v|^2
 \leq(\operatorname{diam}F)^2.
\]
对 \(D\) 中两点取上确界，就得到所需直径估计。

由 \(D\) 的定义，\(R_d(D)=D\)。
若 \(i\in J\)，则 \(i<d\)，并且 \(F_{R_i y}=F_y\)，这里 \(R_i\) 也表示对应的 \(\mathbb R^{d-1}\) 反射。
故 \(\ell(R_i y)=\ell(y)\)，进而 \((y,t)\in D\) 当且仅当 \((R_i y,t)\in D\)。
这说明以前取得的坐标对称性没有被本次操作破坏。
\end{proof}

\begin{proof}[定理\ref{b1:thm:isodiametric}的证明]
空集及直径为零的集合都满足结论。设 \(A\ne\varnothing\)，并记 \(D_0=\operatorname{diam}A>0\)。
若 \(d=1\)，集合包含于 \([\inf A,\sup A]\)，该区间长度为 \(D_0\)，而 \(\omega_1=2\)，故结论成立。

设 \(d\geq2\)。从 \(A_0=A\) 出发，依次对第 \(1,\ldots,d\) 个坐标应用前一引理，得到 \(A_1,\ldots,A_d\)。
它们满足
\[
 m_d^*(A_0)\leq m_d(A_1)\leq\cdots\leq m_d(A_d),
 \quad\operatorname{diam}A_k\leq D_0,
\]
且 \(A_k\) 关于前 \(k\) 个坐标反射不变。
若某一步所得集合为空，上式已经给出 \(m_d^*(A)=0\)。否则对任意 \(x\in A_d\)，连续反射全部坐标得到 \(-x\in A_d\)，故
\[
 2|x|=|x-(-x)|\leq\operatorname{diam}A_d\leq D_0.
\]
因此 \(A_d\subseteq\overline B(0,D_0/2)\)，由体积的伸缩公式可得
\[
 m_d^*(A)\leq m_d(A_d)
 \leq m_d\bigl(\overline B(0,D_0/2)\bigr)
 =\omega_d(D_0/2)^d.
\]
最后，任意半径 \(D_0/2\) 的闭球都具有直径 \(D_0\) 和上述体积，故确实达到等号。
\end{proof}

证明没有把原集合直接装进同直径的球，而是先改变其形状，再比较体积。
中间取闭包可能增大体积，但不增大直径，也不破坏已有的反射对称性；
这正是能够对任意有界集合使用同一论证的原因。
球达到最佳上界，已经足以确定下面的测度常数，不需要再对全部等号情形作分类。

\subsection{\texorpdfstring{\(\mathcal H^d\)}{Hd} 与 Lebesgue 测度}
\label{b1:subsec:150303}

先使用第\ref{b1:sec:1501}节未归一化的覆盖量 \(h_\delta^d\) 与 \(h^d\)。
这样可以先确定几何给出的常数，再与 \(c_d\) 的 Gamma 积分表达式核对。
所用比较见 Fremlin 的《Measure Theory》\cite[264I]{Fremlin2016Measure}。

有一个细节不能省略：Vitali 定理只覆盖几乎所有点，Hausdorff 定义却要求覆盖全部点。
我们先说明 Lebesgue 零测集能用任意小的直径代价补上。
若 \(m_d^*(N)=0\)，给定 \(\delta,\varepsilon>0\)，由\cref{b1:prop:cube-cover}可取半开立方体 \((Q_i)\) 覆盖 \(N\)，使
\[
 \sum_i m_d(Q_i)<\varepsilon d^{-d/2}.
\]
把边长过大的立方体作有限次等分，总代价不变，可使每个新立方体的直径都小于 \(\delta\)。
边长为 \(a\) 的立方体直径为 \(\sqrt d\,a\)，故细分后的覆盖满足
\[
 \sum_i(\operatorname{diam}Q_i)^d
 =d^{d/2}\sum_i m_d(Q_i)<\varepsilon.
\]
这直接从覆盖定义给出 \(h_\delta^d(N)=0\)，进而 \(h^d(N)=0\)，并未预先使用待证的测度一致性。

\begin{proposition}\label{b1:prop:raw-hausdorff-lebesgue}
对任意 \(E\subseteq\mathbb R^d\) 及每个 \(\delta>0\)，都有
\[
 h_\delta^d(E)=h^d(E)=\frac{2^d}{\omega_d}m_d^*(E).
\]
\end{proposition}
\begin{proof}
先证下界。若 \((E_i)\) 是 \(E\) 的一个直径不超过 \(\delta\) 的覆盖，
则由等直径不等式与外测度的次可加性，
\[
 m_d^*(E)\leq\sum_i m_d^*(E_i)
 \leq\frac{\omega_d}{2^d}\sum_i(\operatorname{diam}E_i)^d.
\]
对覆盖取下确界，得到 \(m_d^*(E)\leq(\omega_d/2^d)h_\delta^d(E)\)。
因此若 \(m_d^*(E)=\infty\)，全部所需等式已经成立。

以下设 \(m_d^*(E)<\infty\)。给定 \(\eta>0\)，由\cref{b1:prop:open-cover-approximation}取开集 \(G\supseteq E\)，使
\(m_d(G)<m_d^*(E)+\eta\)。
考虑全部包含于 \(G\)、直径小于 \(\delta\) 的正半径中心闭球。
它们在 \(E\) 上细覆盖，故\cref{b1:thm:vitali-covering-lebesgue}给出至多可数个两两不交的闭球 \((B_j)\)，使
\[
 N=E\setminus\bigcup_jB_j,
 \quad m_d^*(N)=0.
\]
这些球仍包含于 \(G\)，于是
\[
 \sum_j(\operatorname{diam}B_j)^d
 =\frac{2^d}{\omega_d}\sum_j m_d(B_j)
 \leq\frac{2^d}{\omega_d}m_d(G).
\]
再用本小节开头构造的小立方体覆盖 \(N\)，令其直径代价之和小于 \(\varepsilon\)。
合并两族就真正覆盖了 \(E\)，并且每个覆盖集的直径都小于 \(\delta\)。因此
\[
 h_\delta^d(E)
 \leq\frac{2^d}{\omega_d}\bigl(m_d^*(E)+\eta\bigr)+\varepsilon.
\]
先让 \(\varepsilon,\eta\downarrow0\)，得所需上界。
这对每个 \(\delta>0\) 都成立，再取 \(\delta\downarrow0\) 即得到 \(h^d\) 的公式。
\end{proof}

下界使用任意覆盖集的等直径体积界，上界则主动选择球，因为球恰好达到这个界。
Vitali 选择把球的质量之和控制在开集的体积内；剩余零测集由立方体覆盖补齐。
所以两种方向的估计最终给出了同一个常数。

\subsection{归一化常数}
\label{b1:subsec:150304}

第\ref{b1:sec:1501}节对所有实数 \(s\geq0\) 统一采用
\[
 \mathcal H^s=c_sh^s,\quad
 c_s=\frac{\pi^{s/2}}{2^s\Gamma(1+s/2)}.
\]
在整数维数，这个看似解析的表达式恰好等于单位直径球的体积。

\begin{proposition}[整数维归一化]\label{b1:prop:hausdorff-normalization}
对每个正整数 \(d\)，有
\[
 \omega_d\Gamma(1+d/2)=\pi^{d/2},\quad c_d=\frac{\omega_d}{2^d}.
\]
\end{proposition}
\begin{proof}
第\ref{b1:subsec:070504}小节已用二维换元证明
\(\int_{\mathbb R}e^{-t^2}\,\mathrm dt=\sqrt\pi\)。
由 Tonelli 定理反复分离坐标，得到
\[
 \int_{\mathbb R^d}e^{-|x|^2}\,\mathrm dx
 =\left(\int_{\mathbb R}e^{-t^2}\,\mathrm dt\right)^d
 =\pi^{d/2}.
\]
另一方面，恒等式
\(e^{-|x|^2}=\int_{|x|^2}^\infty e^{-u}\,\mathrm du\)
与 Tonelli 定理给出
\[
 \begin{aligned}
 \int_{\mathbb R^d}e^{-|x|^2}\,\mathrm dx
 &=\int_0^\infty e^{-u}
      m_d\bigl(\{\,x:|x|^2\leq u\,\}\bigr)\,\mathrm du\\
 &=\omega_d\int_0^\infty e^{-u}u^{d/2}\,\mathrm du
 =\omega_d\Gamma(1+d/2).
 \end{aligned}
\]
这里仅使用了球体积的伸缩公式，没有使用高维极坐标或球面的面积公式。
比较两次计算，便得第一式，代入 \(c_d\) 的定义得到第二式。
\end{proof}

\begin{theorem}\label{b1:thm:hausdorff-lebesgue}
采用上述归一化时，对任意 \(E\subseteq\mathbb R^d\)，有
\[
 \mathcal H^d(E)=m_d^*(E).
\]
因此 \(\mathcal H^d\) 与 Lebesgue 测度具有相同的 Carathéodory 可测集，
且在这些集合上完全相等。
\end{theorem}
\begin{proof}
由前一命题及\cref{b1:prop:raw-hausdorff-lebesgue}，
\[
 \mathcal H^d(E)
 =c_dh^d(E)
 =\frac{\omega_d}{2^d}\cdot\frac{2^d}{\omega_d}m_d^*(E)
 =m_d^*(E).
\]
这首先是两个外测度在所有集合上的相等。
将其代入 Carathéodory 可测性判据，可知可测集族也相同，限制后的测度自然一致。
\end{proof}

\begin{corollary}\label{b1:cor:hausdorff-flat-subspace}
设 \(1\leq k\leq d\)，\(T\colon\mathbb R^k\to\mathbb R^d\) 是等距仿射嵌入。
则对任意 \(E\subseteq\mathbb R^k\)，有
\[
 \mathcal H^k_{\mathbb R^d}\bigl(T(E)\bigr)=m_k^*(E).
\]
这里的下标仅用于区分计算 Hausdorff 测度的环境空间。
\end{corollary}
\begin{proof}
映射 \(T\) 及其从像空间到 \(\mathbb R^k\) 的逆映射都保持距离。
由\cref{b1:prop:hausdorff-lipschitz}中的等距不变性与子空间环境独立性，
\[
 \mathcal H^k_{\mathbb R^d}\bigl(T(E)\bigr)
 =\mathcal H^k_{\mathbb R^k}(E).
\]
再在 \(k\) 维欧氏空间应用前一定理，得到右侧等于 \(m_k^*(E)\)。
\end{proof}

因此一条直线上的 \(\mathcal H^1\) 是通常的长度，一个平面上的 \(\mathcal H^2\) 是通常的面积，
不论它们嵌入多高维的欧氏空间。
这些是等距参数化的直接结果；曲面或非等距参数化中的体积变化，需要另外计算。

例如 \(c_1=1\)、\(c_2=\pi/4\)、\(c_3=\pi/6\)，分别把直径代价换算成线段长度、圆盘面积和球体积。
这些数值也可由 Gamma 积分直接核对：分部积分给出
\(\Gamma(t+1)=t\Gamma(t)\)（\(t>0\)），而 \(\Gamma(1)=1\)、
\(\Gamma(1/2)=2\int_0^\infty e^{-v^2}\,\mathrm dv=\sqrt\pi\)。
分部积分先在 \([\varepsilon,R]\) 上进行，再令 \(\varepsilon\downarrow0\)、\(R\uparrow\infty\)，
端点项因 \(t>0\) 与指数衰减而消失。
零维时 \(c_0=1\)，仍得到第\ref{b1:sec:1501}节的计数测度。

归一化不会改变某个集合的 Hausdorff 维数，因为对固定 \(s\)，因子 \(c_s\) 是严格正的有限数。
但涉及长度、面积、体积或密度中的精确常数时，归一化必须保持一致。
比较不同参考书的公式，应同时检查覆盖代价究竟写作 \((\operatorname{diam}E_i)^s\)，
还是已经乘入 \(c_s\)，不能只根据符号 \(\mathcal H^s\) 判断它们数值相同。

% !TeX root = ../../Book1.tex
\section{Frostman 引理}
\label{b1:sec:1504}

在 Cantor 集上，基本区间已经告诉我们应把质量放在哪里。
一般紧集没有现成的分支结构，却仍可以提出同一要求：半径为 \(r\) 的球至多承载常数倍的 \(r^s\) 质量。
质量分布原理说明，这种要求足以给出维数下界。
\term[frostman-yinli]{Frostman 引理}[Frostman's lemma]则回答反方向的问题：集合的 Hausdorff 测度为正时，是否总能找到这样的质量分布？
本节对欧氏空间中的紧集给出肯定回答，并在本节内证明构造所需的紧集概率测度序列紧性。

% 来源范围：Fremlin2016Measure 471Xd(ii)只作为Frostman型结论的对照，
% 不把该习题的附加假设误当成一般紧集定理的完整证明。
% 下文有限树递推、割集下界、原子测度弱极限及能量估计均在本书内完整展开。
% Mattila1995Geometry与Falconer2014Fractal在本节只作专题伴读推荐。

\subsection{质量分布原理回顾与统一增长条件}
\label{b1:subsec:150401}

\cref{b1:lem:mass-distribution}已经正式给出了质量分布原理及其证明。
本小节只回顾其中的覆盖机制，并强调 Frostman 理论所需估计必须在中心、常数与尺度上统一。
设 \(s>0\)，\(\mu\) 是集中在 Borel 集 \(E\subset\mathbb R^d\) 上的有限正测度，
\(\mu(E)>0\)。若所有足够小的球都满足
\[
 \mu\bigl(B(x,r)\bigr)\leq A r^s,
\]
则任意细覆盖都必须付出一定的直径幂和。
事实上，对与 \(E\) 相交且直径为 \(a>0\) 的覆盖集 \(U\)，取 \(x\in U\cap E\)，
就有 \(U\subset\overline B(x,a)\)，从而其质量至多为 \(Aa^s\)。
这里不要求 \(U\) 可测：应写作 \(\mu^*(U\cap E)\)，再用包住它的可测球估计。
这正是\cref{b1:lem:mass-distribution}的覆盖论证。

值得注意的是不等式的方向。若希望证明集合至少有 \(s\) 维，
就应限制小球不能装入过多的质量，而不是要求小球质量至少为 \(r^s\)。
一旦总质量固定，单个小集合能承载的质量越少，覆盖整个集合所需的集合就越多。
常数 \(A\) 只改变最后的测度下界，不改变维数。

若估计只在 \(\mu\)-几乎处处中心成立，也常可先限制到一个正质量子集上。
但必须检查两个统一性问题：常数是否统一，以及允许的半径范围是否统一。
例如，若对每个 \(x\) 都有某个 \(A_x<\infty\) 和 \(r_x>0\)，
可按 \(A_x\leq n\)、\(r_x\geq1/n\) 分组；在分组可测的条件下，
总有一组仍有正质量。仅写“对每个点有局部估计”而不作这一步，不能直接套用质量分布原理。

\subsection{Frostman 测度}
\label{b1:subsec:150402}

\begin{definition}\label{b1:def:frostman-measure}
设 \(s\geq0\)，\(K\subset\mathbb R^d\) 为非空紧集。
若非零有限正 Borel 测度 \(\mu\) 集中在 \(K\) 上，并存在 \(A<\infty\)，使
\[
 \mu\bigl(B(x,r)\bigr)\leq A r^s
 \quad(x\in\mathbb R^d,\ r>0),
\]
则称 \(\mu\) 为 \(K\) 上的一个 \term[frostman-cedu]{Frostman 测度}[Frostman measure]，
其中 \(s\) 是所用的增长指数。
\end{definition}

“集中在 \(K\) 上”只要求 \(\mu(\mathbb R^d\setminus K)=0\)，
不要求 \(\operatorname{supp}\mu=K\)。
由于总质量有限，除以 \(\mu(K)\) 即可改成概率测度，只需相应调整 \(A\)。
当 \(s>0\) 时，这样的测度没有原子，因为令 \(r\downarrow0\) 就得 \(\mu(\{x\})=0\)。
当 \(s=0\) 时，非空紧集上的任意 Dirac 测度均满足要求。

\begin{lemma}\label{b1:lem:hausdorff-content-positive}
若 \(s>0\)，则对任意集合 \(E\subset\mathbb R^d\)，
\[
 \mathcal H^s(E)>0\quad\Longleftrightarrow\quad
 \mathcal H^s_\infty(E)>0.
\]
\end{lemma}
\begin{proof}
这正是\cref{b1:prop:hausdorff-borel-hull}中零测性等价关系的逆否形式：
无尺度限制的代价能够任意小时，各个覆盖集的直径也被迫任意小。
\end{proof}

下面的构造不把 \(\mathcal H^s|_K\) 归一化，因为 \(\mathcal H^s(K)\) 可能为无穷。
我们先在有限层的二进立方体上分配质量，再把层数趋于无穷。

这里的二进立方体是半开立方体 \(Q_0\) 及反复等分各边所得的子立方体；
每个立方体有 \(2^d\) 个互不相交的子立方体。记其边长为 \(\ell(Q)\)。
半开约定使同层立方体真正构成分割，点不会因落在公共边界而被重复计数。

\cref{b1:fig:frostman-dyadic-tree}示意有限树上的两个约束：质量从根向下守恒，
每个结点所得质量又不超过递推给出的容量 \(F_N(Q)\)。末层质量最终放在
\(K\) 中选定的点上，形成离散测度。

\begin{figure}[htbp]
\centering
\begin{tikzpicture}[>=stealth,line cap=round,line join=round,
  box/.style={draw,rounded corners,align=center,inner sep=4pt}]
  \node[box] (root) at (0,2.2) {$Q_0$\\$m(Q_0)=a$};
  \node[box] (left) at (-3.4,0.5) {$Q'$\\$m(Q')\leq F_N(Q')$};
  \node at (0,0.5) {$\cdots$};
  \node[box] (right) at (3.4,0.5) {$Q''$\\$m(Q'')\leq F_N(Q'')$};
  \node[box] (ll) at (-4.5,-1.35) {$R_1$\\$m(R_1)\delta_{x_{R_1}}$};
  \node[box] (lr) at (-2.3,-1.35) {$R_2$\\$m(R_2)\delta_{x_{R_2}}$};
  \node at (2.25,-1.35) {$\cdots$};
  \node[box] (rr) at (4.15,-1.35) {$R$\\$m(R)\delta_{x_R}$};
  \draw[->,thick] (root)--(left);
  \draw[->,thick] (root)--(right);
  \draw[->,thick] (left)--(ll);
  \draw[->,thick] (left)--(lr);
  \draw[->,thick] (right)--(rr);
  \node[align=center] at (0,-1.35) {逐层分配\\质量守恒};
\end{tikzpicture}
\caption[二进质量分配]{二进质量分配。根质量 \(a\) 逐层向子立方体分配；每个结点受容量 \(F_N\) 控制，末层质量则变成位于 \(K\) 上的原子。图中只画出部分分支。}
\label{b1:fig:frostman-dyadic-tree}
\end{figure}

下面把构造中唯一需要的紧性输入单独列出。这样本节不依赖\chapref{b1:ch:12}的完整弱收敛理论；
该章给出的是波兰空间、紧性与 Portmanteau 判据的更一般版本。

\begin{lemma}[紧集上概率测度的序列紧性]\label{b1:lem:compact-probability-subsequence}
设 \(K\subset\mathbb R^d\) 为非空紧集，\((\nu_n)\) 是 \(K\) 上的 Borel 概率测度序列。
则存在子列 \((\nu_{n_j})\) 与 \(K\) 上的 Borel 概率测度 \(\nu\)，使
\[
 \int_K f\,\mathrm d\nu_{n_j}\longrightarrow\int_K f\,\mathrm d\nu
 \quad\bigl(f\in C(K;\mathbb R)\bigr).
\]
此外，对每个在 \(K\) 中相对开的集合 \(G\)，都有
\[
 \nu(G)\leq\liminf_{j\to\infty}\nu_{n_j}(G).
\]
\end{lemma}
\begin{proof}
先说明 \(C(K)\) 可分。对每个 \(n\geq1\)，取有限的 \(1/n\)-网
\(\{q_{n,1},\ldots,q_{n,N_n}\}\)，并令
\[
 w_{n,k}(x)=\max\{0,2/n-|x-q_{n,k}|\},
 \qquad
 \phi_{n,k}(x)=\frac{w_{n,k}(x)}{\sum_iw_{n,i}(x)}.
\]
分母处处为正，且 \(\sum_k\phi_{n,k}=1\)。连续函数在紧集上一致连续，故
\(\sum_k f(q_{n,k})\phi_{n,k}\) 随 \(n\to\infty\) 一致逼近 \(f\)。
再把有限个系数换成有理数，就得到 \(C(K;\mathbb R)\) 的可数稠密族。

由于 \(K\) 紧，\(C(K;\mathbb R)=C_0(K;\mathbb R)\)。由\cref{b1:cor:positive-c0-representation}，
概率测度恰对应于 \(C(K;\mathbb R)^*\) 中满足
\[
 L(1)=1,\qquad L(f)\geq0\quad(f\geq0)
\]
的泛函；它们都位于对偶单位球中。这些条件在弱-*拓扑下闭。
\cref{b1:thm:banach-alaoglu}、\cref{b1:thm:weak-star-metrizable}与 \(C(K)\) 的可分性说明，
这族泛函是紧致可度量空间，因而每个序列都有弱-*收敛子列。
再用同一个表示定理识别极限泛函，就得到概率测度 \(\nu\) 及题设的积分收敛。

最后设 \(G\subset K\) 相对开。若 \(G=K\)，所需不等式由总质量为 \(1\) 得到；
否则令
\[
 f_m(x)=\min\{1,m\,\operatorname{dist}(x,K\setminus G)\}.
\]
则 \(f_m\in C(K)\)、\(0\leq f_m\leq\mathbf1_G\)，且 \(f_m\uparrow\mathbf1_G\)。
于是对每个 \(m\)，
\[
 \int_K f_m\,\mathrm d\nu
 =\lim_{j\to\infty}\int_K f_m\,\mathrm d\nu_{n_j}
 \leq\liminf_{j\to\infty}\nu_{n_j}(G).
\]
令 \(m\to\infty\)，再用单调收敛定理，即得开集不等式。
\end{proof}

\begin{theorem}[Frostman]\label{b1:thm:frostman}
设 \(K\subset\mathbb R^d\) 为非空紧集，\(0<s\leq d\)。
则 \(\mathcal H^s(K)>0\) 当且仅当 \(K\) 上存在增长指数为 \(s\) 的 Frostman 测度。
\end{theorem}
存在性方向有四个转换：先把正覆盖量变成有限二进立方体树上的容量下界；
再从根向下分配质量，得到具有统一立方体控制的离散测度；
随后用前一引理抽取弱收敛子列；最后用开集不等式把球质量控制传给极限。即
\[
 \begin{aligned}
 \mathcal H^s_\infty(K)>0
 &\longrightarrow \text{有限树容量与质量分配}
 \longrightarrow \mu_N(B(x,r))\lesssim r^s\\
 &\longrightarrow \mu_{N_j}\Rightarrow\mu
 \longrightarrow \mu(B(x,r))\lesssim r^s.
 \end{aligned}
\]
\begin{proof}
已有测度时，由 \(\overline B(x,r)\subset B(x,2r)\) 及质量分布原理即得
\(\mathcal H^s(K)>0\)。
以下设 \(\mathcal H^s(K)>0\)，取内部包含 \(K\) 的半开立方体 \(Q_0\)，边长为 \(L\)。
由前一引理，
\[
 a=\frac{\mathcal H^s_\infty(K)}{c_s d^{s/2}}>0.
\]
又因 \(Q_0\) 本身覆盖 \(K\)，有 \(a\leq L^s\)。

固定层数 \(N\)。给第 \(N\) 层立方体规定数值
\[
 F_N(Q)=
 \begin{cases}
  \ell(Q)^s,&Q\cap K\ne\varnothing,\\
  0,&Q\cap K=\varnothing.
 \end{cases}
\]
然后逐层向上定义
\begin{equation}\label{b1:eq:frostman-tree-recursion}
 F_N(Q)=\min\left\{\ell(Q)^s,\ \sum_{Q'\text{ 为 }Q\text{ 的子立方体}}F_N(Q')\right\}.
\end{equation}
这个递推有一个直接的覆盖解释：在 \(Q\) 内，要覆盖所有与 \(K\) 相交的末层立方体，
可以选取 \(Q\) 本身，也可以分别在各个子立方体内继续选取。
对有限层数归纳，\(F_N(Q_0)\) 就是这类选取所能达到的最小边长幂和。
更明确地说，每条通往非空末层立方体的分支至少须选中一个立方体；
若选中了某立方体，就不必再选它的后代。
每一种如此得到的有限族都覆盖 \(K\)，所以
\[
 c_s d^{s/2} F_N(Q_0)\geq\mathcal H^s_\infty(K),
 \quad F_N(Q_0)\geq a.
\]

现在从上向下分配质量。给根立方体质量 \(a\)；
若 \(Q\) 已分得 \(m(Q)\leq F_N(Q)\)，由%
\eqref{b1:eq:frostman-tree-recursion}可把它分成非负数 \(m(Q')\leq F_N(Q')\)，
且所有子立方体所得质量之和恰为 \(m(Q)\)。
这种分配总能完成：依次向子立方体分配不超过其限额的质量，
由于所有限额之和至少为 \(m(Q)\)，不会有剩余质量无处安放。
有限次重复后，在每个与 \(K\) 相交的末层立方体内选一点 \(x_Q\)，令
\[
 \mu_N=\sum_{Q\text{ 在第 }N\text{ 层}}m(Q)\delta_{x_Q},
\]
空立方体不计入和式。于是 \(\mu_N(K)=a\)，
而每个不深于第 \(N\) 层的立方体均满足
\[
 \mu_N(Q)=m(Q)\leq\ell(Q)^s.
\]
这里使用了半开分割，故原子即使位于某些立方体边界上，等式仍然成立。

概率测度 \(\mu_N/a\) 全部集中在紧集 \(K\) 上。
由\cref{b1:lem:compact-probability-subsequence}，可取子列 \((\mu_{N_j}/a)\)
弱收敛到 \(K\) 上的概率测度；把极限乘以 \(a\)，并视为
\(\mathbb R^d\) 上的测度，记为 \(\mu\)。
因此 \(\mu(K)=a>0\)。还须证明小球质量估计在极限中保留；
不能直接对半开立方体使用闭集版本的 Portmanteau 不等式。

若 \(0<r<L\)，选取一层，使其边长 \(\ell\) 满足 \(r\leq\ell<2r\)。
球 \(B(x,r)\) 在每个坐标方向至多遇到四个该层区间，
故它至多与 \(4^d\) 个该层立方体相交。只要 \(N\) 不小于这一层数，就有
\[
 \mu_N\bigl(B(x,r)\bigr)\leq4^d\ell^s\leq4^d2^s r^s.
\]
对 \(K\cap B(x,r)\) 使用\cref{b1:lem:compact-probability-subsequence}中的开集不等式，得
\[
 \mu\bigl(B(x,r)\bigr)
 \leq\liminf_{j\to\infty}\mu_{N_j}\bigl(B(x,r)\bigr)
 \leq4^d2^s r^s,
\]
其中只需取 \(j\) 充分大，使 \(N_j\) 不小于所选层数。
若 \(r\geq L\)，则 \(\mu\bigl(B(x,r)\bigr)\leq a\leq L^s\leq r^s\)。
故 \(\mu\) 满足定义中的统一增长估计。
\end{proof}

紧性在这段证明中有明确的位置：它使各层构造的质量无法在弱极限中逃走。
另一方面，各层的质量分配不必彼此相容；相容性由取弱极限解决。
把这两件事分开，往往比直接在一棵无限树上规定所有分支的质量更方便。
关于更一般集合上的版本及其与势论的关系，可继续阅读 Mattila 的
《Geometry of Sets and Measures in Euclidean Spaces》\cite{Mattila1995Geometry}；
这里已把所需的紧集概率测度序列紧性和开集不等式在本节内证明；
\chapref{b1:ch:12}则把这些结论放到更一般的测度弱收敛框架中。

\subsection{Hausdorff 维数下界}
\label{b1:subsec:150403}

\begin{corollary}\label{b1:cor:frostman-dimension}
对非空紧集 \(K\subset\mathbb R^d\)，
\[
 \dim_{\mathrm H}K=
 \sup\{\,s\geq0\mid K\text{ 上存在增长指数为 }s\text{ 的 Frostman 测度}\,\}.
\]
\end{corollary}
\begin{proof}
若存在增长指数为 \(s\) 的测度，质量分布原理给出 \(\dim_{\mathrm H}K\geq s\)。
反之，若 \(0<s<\dim_{\mathrm H}K\)，则 \(\mathcal H^s(K)=\infty\)，
由 Frostman 引理可构造所需测度。
当 \(\dim_{\mathrm H}K=0\) 时，Dirac 测度给出指数 \(0\)，也得到等式。
\end{proof}

这不是说临界指数一定可以取到。设 \(s_0=\dim_{\mathrm H}K>0\)，
在 \(s=s_0\) 时，能否找到相应测度，恰好取决于
\(\mathcal H^{s_0}(K)\) 是否为正；维数本身并不决定这个数值。
在三分 Cantor 集上，前面构造的等分质量恰好实现临界指数。
对于别的集合，先证明所有 \(s<s_0\) 都有增长估计，已经足以得到维数下界，
不应为获得一个并非必要的端点估计而停下整个论证。

作为应用，设 \(F:K\rightarrow\mathbb R^m\) 是双 Lipschitz 映射，
\(\mu\) 是 \(K\) 上的一个 \(s\) 指数 Frostman 测度。
若逆映射的 Lipschitz 常数为 \(L\)，并且 \(B(y,r)\cap F(K)\ne\varnothing\)，
任选其中一点 \(F(x_0)\)，便有
\[
 F^{-1}\bigl(B(y,r)\cap F(K)\bigr)\subset B(x_0,2Lr).
\]
所以推前测度 \(F_\#\mu\) 仍满足 \(s\) 次增长估计。
这从测度构造一侧再次解释了双 Lipschitz 不变性：
距离只改变固定倍数时，质量增长的指数不会改变。

\subsection{能量方法初步}
\label{b1:subsec:150404}

逐球估计要求控制所有中心和尺度，有时不便直接验证。
另一种办法是把点对之间的距离作一次非负积分。

\begin{definition}\label{b1:def:riesz-energy}
设 \(t>0\)，\(\mu\) 为 \(\mathbb R^d\) 上的有限正 Borel 测度。
定义其 \term[riesz-nengliang]{Riesz 能量}[Riesz energy] 为
\[
 I_t(\mu)\coloneqq\iint_{\mathbb R^d\times\mathbb R^d}|x-y|^{-t}\,d\mu(y)\,d\mu(x),
\]
并约定对角线上的核值为 \(+\infty\)。
相应的非负函数
\[
 U_t^\mu(x)\coloneqq\int_{\mathbb R^d}|x-y|^{-t}\,d\mu(y)
\]
称为这里的 \(t\) 次势。能量与势都允许取无穷值。
\end{definition}
\symidx{\(I_t(\mu)\)：Riesz 能量}
\symidx{\(U_t^\mu\)：距离负幂核的势}

核是非负 Borel 函数，故 Tonelli 定理保证上述积分有意义，
也保证 \(U_t^\mu\) 可测。这里用距离负幂 \(t\) 作为下标；
其他书中若按势算子的阶数编号，可能使用核 \(|x-y|^{\alpha-d}\)，
其参数应对应为 \(t=d-\alpha\)，不能只按字母名称比较。

\begin{proposition}\label{b1:prop:frostman-finite-energy}
若非零有限正测度 \(\mu\) 满足
\(\mu\bigl(B(x,r)\bigr)\leq A r^s\) 对所有 \(x,r\) 成立，
则对 \(0<t<s\) 有 \(I_t(\mu)<\infty\)。
\end{proposition}
\begin{proof}
记 \(M=\mu(\mathbb R^d)\)。由
\[
 u^{-t}=\int_u^\infty t r^{-t-1}\,dr
\]
及 Tonelli 定理，包括 \(u=0\) 的无穷值情形，得到
\begin{equation}\label{b1:eq:energy-ball-integral}
 I_t(\mu)=t\int_0^\infty r^{-t-1}
       \int_{\mathbb R^d}\mu\bigl(B(x,r)\bigr)\,d\mu(x)\,dr.
\end{equation}
在 \(0<r<1\) 用增长估计，在 \(r\geq1\) 用总质量估计，即得
\[
 I_t(\mu)\leq tAM\int_0^1r^{s-t-1}\,dr+
 tM^2\int_1^\infty r^{-t-1}\,dr
 =\frac{tAM}{s-t}+M^2<\infty.
\]
\end{proof}

\begin{proposition}\label{b1:prop:finite-energy-dimension}
设非零有限正测度 \(\mu\) 集中在 Borel 集 \(E\) 上。
若 \(I_t(\mu)<\infty\)，则 \(\dim_{\mathrm H}E\geq t\)。
\end{proposition}
\begin{proof}
由能量有限，\(U_t^\mu(x)<\infty\) 对 \(\mu\)-几乎处处的 \(x\) 成立。
可选正整数 \(n\)，使 Borel 集
\[
 E_n=\{\,x\in E\mid U_t^\mu(x)\leq n\,\}
\]
有正质量，令 \(\nu=\mu|_{E_n}\)。
若 \(B(z,r)\cap E_n\ne\varnothing\)，取 \(x\) 属于此交集。
对 \(y\in B(z,r)\) 有 \(|x-y|<2r\)，故
\[
 \nu\bigl(B(z,r)\bigr)
 \leq(2r)^t\int_{B(z,r)}|x-y|^{-t}\,d\mu(y)
 \leq 2^t n r^t.
\]
若交集为空，左边为零，同一估计仍成立。
对正质量测度 \(\nu\) 使用质量分布原理，便得
\(\mathcal H^t(E_n)>0\)，因而 \(\dim_{\mathrm H}E\geq t\)。
\end{proof}

有限能量先给出势几乎处处有限，再限制到势有统一上界的部分；
这个限制步骤不能省略为“所以原测度具有统一增长估计”。
而且有限能量必定排除原子：若某点质量为 \(b>0\)，
对角线上的该点对就有质量 \(b^2\)，使能量为无穷。

结合以上两命题与 Frostman 引理，对非空紧集 \(K\) 有
\[
 \dim_{\mathrm H}K=
 \sup\Bigl(\{0\}\cup
 \{\,t>0\mid \text{存在集中在 }K\text{ 上的概率测度 }\mu,\ I_t(\mu)<\infty\,\}\Bigr).
\]
其中下界使用有限能量命题；上界的反向逼近则对
\(0<t<s<\dim_{\mathrm H}K\) 先构造 \(s\) 指数 Frostman 测度。
公式中的严格不等式 \(t<s\) 有实质作用，在端点把积分
\(\int_0^1r^{-1}\,dr\) 当成有限会得出错误结论。
Falconer 的《Fractal Geometry: Mathematical Foundations and Applications》%
\cite{Falconer2014Fractal}可作为继续学习维数估计与分形应用的伴读；
本节只建立了距离负幂核所需的基本积分估计，尚未引入一般容量理论。

% !TeX root = ../../Book1.tex
\section{Hausdorff 密度}
\label{b1:sec:1505}

维数给出的是尺度变化的指数，尚未描述质量在某一点附近如何分布。
即使两个集合的维数相同，它们在小球中的质量也可能有不同的渐近行为。
本节把上一章的密度问题改写为以 \(r^s\) 为分母的几何问题。
整数维平面提供归一化的基准；一般集合则先有上密度的几乎处处估计，
不能预先假定密度极限存在。

% 实际读取：Simon1983Lectures §3，印刷页10--17，密度定义及3.2、3.6。
% 原书的一般度量空间版本与本节不同；本节限R^d的Borel E且H^s(E)<infty，
% 上界改用本书13.2的Radon细覆盖，下界完整写出统一半径分组。
% 平面密度及Cantor端点的两列尺度计算在本文直接推导。

\subsection{\texorpdfstring{\(s\)}{s}-维密度}
\label{b1:subsec:150501}

沿用 \(c_s\) 的约定，记
\[
 \alpha_s=2^s c_s=\frac{\pi^{s/2}}{\Gamma(1+s/2)}.
\]
对正整数 \(k\)，第\ref{b1:sec:1503}节证明了 \(\alpha_k=\omega_k\)，
即 \(\mathbb R^k\) 中单位球的体积。
非整数 \(s\) 时，\(\alpha_s\) 只是明确指定的正归一化常数，
并不意味着存在一个通常意义的“\(s\) 维欧氏空间”。

\begin{definition}\label{b1:def:hausdorff-density}
设 \(s>0\)，\(\mu\) 为 \(\mathbb R^d\) 上的正 Radon 测度。
定义 \(\mu\) 在 \(x\) 处的 \term[shang-hausdorff-midu]{上 \(s\) 维密度}[upper \(s\)-dimensional density]
和 \term[xia-hausdorff-midu]{下 \(s\) 维密度}[lower \(s\)-dimensional density] 为
\[
 \Theta^{*s}(\mu,x)=
 \limsup_{r\downarrow0}\frac{\mu\bigl(\overline B(x,r)\bigr)}{\alpha_s r^s},
 \quad
 \Theta_*^s(\mu,x)=
 \liminf_{r\downarrow0}\frac{\mu\bigl(\overline B(x,r)\bigr)}{\alpha_s r^s}.
\]
若二者相等，则把共同值记为 \(\Theta^s(\mu,x)\)，称为 \(s\) 维密度。
对 Borel 集 \(E\) 及局部有限测度 \(\mathcal H^s|_E\)，还记
\[
 \Theta^{*s}(E,x)=\Theta^{*s}(\mathcal H^s|_E,x),
 \quad
 \Theta_*^s(E,x)=\Theta_*^s(\mathcal H^s|_E,x).
\]
\end{definition}
\symidx{\(\Theta^{*s},\Theta_*^s,\Theta^s\)：上、下与存在时的 \(s\) 维密度}

这些函数是 Borel 可测的。事实上，固定 \(r\) 时的闭球质量为上半连续函数；
利用闭球从外逼近的连续性，可把半径上、下极限化为可数有理半径的上、下极限。
论证与\cref{b1:prop:relative-density-borel}相同，这里的分母
\(\alpha_s r^s\) 更是处处正的连续函数。
同样，把闭球改成开球不改变这些上、下极限：
固定 \(0<\varepsilon<1\)，使用
\[
 B(x,r)\subset\overline B(x,r)\subset B\bigl(x,(1+\varepsilon)r\bigr)
\]
夹住比值，先取上、下极限，再令 \(\varepsilon\downarrow0\)。
这个理由不需要任何球面都是零测集。

若 \(\mu\) 是 \(\mathbb R^d\) 中的一个原子测度，并且 \(\mu(\{x\})>0\)，
则对每个 \(s>0\) 都有 \(\Theta^s(\mu,x)=\infty\)。
相反，一条直线上的长度测度在直线外的点附近为零。
由此可见，密度既依赖增长指数，也依赖观察点与质量支撑的关系。

\subsection{密度存在问题}
\label{b1:subsec:150502}

不能把\cref{b1:thm:radon-measure-differentiation}直接解释为本节密度几乎处处存在。
该定理的分母是一个固定 Radon 测度的小球质量，
而本节的 \(r^s\) 在非整数维情形并没有相应的环境体积测度。
即使 \(s=k<d\) 为整数，\(\mathcal H^k\) 在环境空间 \(\mathbb R^d\) 上也不是局部有限测度。
正确的出发点是下面的估计，它只涉及上密度。

\begin{theorem}[Hausdorff 上密度估计]\label{b1:thm:hausdorff-upper-density}
设 \(E\subset\mathbb R^d\) 为 Borel 集，\(s>0\)，\(\mathcal H^s(E)<\infty\)。
则对 \(\mathcal H^s\)-几乎处处的 \(x\in E\)，
\[
 2^{-s}\leq\Theta^{*s}(E,x)\leq1.
\]
\end{theorem}
\begin{proof}
令 \(\mu=\mathcal H^s|_E\)。它是有限 Borel 测度，
由欧氏空间中有限 Borel 测度的正则性，亦为 Radon 测度。
若 \(\mu(E)=0\)，结论自动成立，以下允许 \(\mu(E)>0\)。

先证上界。固定 \(t>1\)，记
\[
 A=\{\,x\in E\mid\Theta^{*s}(E,x)>t\,\}.
\]
只需证明 \(A\) 的每个有界 Borel 部分 \(A_0\) 都为零测集。
给定 \(\varepsilon>0\)，取开集 \(U\supset A_0\)，使
\(\mu(U)\leq\mu(A_0)+\varepsilon\)。
对任意 \(\delta>0\)，由上极限的定义，
每个 \(x\in A_0\) 都是任意小闭球的中心，这些闭球满足
\[
 \overline B(x,r)\subset U,\quad 2r<\delta,\quad
 \mu\bigl(\overline B(x,r)\bigr)>t\alpha_s r^s.
\]
由\cref{b1:cor:radon-vitali-covering}，可选两两不交的至多可数个这样的球
\(\overline B(x_j,r_j)\)，覆盖 \(A_0\) 除去一个 \(\mu\)-零测集。
因遗漏集包含于 \(E\)，它也是 \(\mathcal H^s\)-零测集；
这就允许在 Hausdorff 覆盖估计中真正忽略遗漏部分。因此
\[
 \mathcal H^s_\delta(A_0)
 \leq\sum_j c_s(2r_j)^s
 \leq\frac1t\sum_j\mu\bigl(\overline B(x_j,r_j)\bigr)
 \leq\frac{\mu(A_0)+\varepsilon}{t}.
\]
依次令 \(\delta\downarrow0\)、\(\varepsilon\downarrow0\)，并用
\(\mathcal H^s(A_0)=\mu(A_0)<\infty\)，得到
\(\mu(A_0)\leq\mu(A_0)/t\)，故 \(\mu(A_0)=0\)。
用有界球穷竭 \(A\)，再令 \(t\) 沿可数列下降到 \(1\)，即得上界。

再证下界。固定 \(0<t<2^{-s}\)，令
\[
 A_n=\left\{\,x\in E\ \middle|\
       \mu\bigl(\overline B(x,r)\bigr)\leq t\alpha_s r^s
       \text{ 对所有 }0<r<1/n\,\right\}.
\]
这些集是 Borel 集：只需检查有理 \(r\)，再用闭球从外逼近。
若 \(\Theta^{*s}(E,x)<t\)，则 \(x\) 属于某个 \(A_n\)。
取 \(\delta<1/n\) 及 \(A_n\) 的任意可数 \(\delta\)-覆盖 \((U_j)\)，
可删去不与 \(A_n\) 相交的覆盖集。
对直径 \(a_j>0\) 的 \(U_j\)，取 \(x_j\in U_j\cap A_n\)，得
\[
 \mu^*(U_j\cap A_n)
 \leq\mu\bigl(\overline B(x_j,a_j)\bigr)
 \leq t\alpha_s a_j^s.
\]
直径为零的覆盖集至多为单点，因 \(s>0\)，其 \(\mathcal H^s\) 测度与 \(\mu\) 质量均为零。
于是
\[
 \mu(A_n)\leq t2^s c_s\sum_j a_j^s.
\]
先对覆盖取下确界，再令 \(\delta\downarrow0\)，得到
\[
 \mu(A_n)\leq t2^s\mathcal H^s(A_n)=t2^s\mu(A_n).
\]
由 \(t2^s<1\) 及 \(\mu(A_n)<\infty\)，知 \(\mu(A_n)=0\)。
最后令 \(t\) 沿可数列上升到 \(2^{-s}\)，便排除了上密度小于 \(2^{-s}\) 的点集。
\end{proof}

上界使用可选取的细球覆盖，下界使用任意细覆盖；
常数 \(2^{-s}\) 来自把一个直径为 \(a\) 的任意集合放入半径为 \(a\) 的球。
这也解释了为何不能在下界中擅自用半径 \(a/2\)：等直径不等式比较体积，
并没有断言每个直径为 \(a\) 的集合都包含于某个半径为 \(a/2\) 的球。
Simon 的《Lectures on Geometric Measure Theory》第3节%
\cite{Simon1983Lectures}讨论了这类密度与覆盖的联系；
上面的证明利用本书前两章已有的 Radon 细覆盖，把所需有限性与可测性条件单独列明。

若只知 \(\mathcal H^s|_E\) 局部有限，同一结论仍可局部使用。
例如对 \(E\cap B(0,n)\) 应用定理，在内部点的小球上，该限制与原集合给出相同的质量。
再作可数穷竭即可。这里仍不能删去局部有限性而直接把整个
\(\mathcal H^s\) 当作环境中的 Radon 测度。

\subsection{局部结构}
\label{b1:subsec:150503}

先看平面上的密度。设 \(P\) 为 \(\mathbb R^d\) 中的 \(k\) 维仿射平面，
\(1\leq k\leq d\)，\(T:\mathbb R^k\rightarrow P\) 为等距参数化。
对任意 \(A\subset P\)，
\begin{equation}\label{b1:eq:hausdorff-plane-identification}
 \mathcal H^k_{\mathbb R^d}(A)
 =\mathcal H^k_{\mathbb R^k}\bigl(T^{-1}(A)\bigr)
 =m_k^*\bigl(T^{-1}(A)\bigr).
\end{equation}
第一个等号需注意覆盖所在的空间：一方面，平面内覆盖也是环境中的覆盖；
另一方面，把环境中的覆盖集正交投影到 \(P\)，直径不会增加。
再用等距参数化即可得到反向不等式。
第二个等号是\cref{b1:thm:hausdorff-lebesgue}。

\begin{proposition}\label{b1:prop:plane-hausdorff-density}
若 \(E\subset P\) 为 Borel 集，\(\mathcal H^k|_E\) 局部有限，
则对 \(\mathcal H^k\)-几乎处处的 \(x\in E\)，
\[
 \Theta^k(E,x)=1.
\]
\end{proposition}
\begin{proof}
把 \(E\) 用 \(T^{-1}\) 搬回 \(\mathbb R^k\)。
对 \(x=T(u)\)，球 \(\overline B(x,r)\) 与 \(P\) 的交集对应于
\(\mathbb R^k\) 中的闭球 \(\overline B(u,r)\)，体积为 \(\omega_k r^k\)。
因此所需比值正是集合 \(T^{-1}(E)\) 的 Lebesgue 密度比值。
由\cref{b1:thm:lebesgue-density}及%
\eqref{b1:eq:hausdorff-plane-identification}即得结论。
\end{proof}

这个结论不要求 \(E\) 在平面内开，也不要求它的边界光滑。
更一般地，若 \(g\geq0\) 在平面上局部可积，
\(\nu=g\,\mathcal H^k|_P\)，则由 \(k\) 维 Lebesgue 微分定理，
\[
 \Theta^k(\nu,x)=g(x)
 \quad\text{对 }\mathcal H^k\text{-几乎处处的 }x\in P.
\]
归一化后，一个平面上均匀的单位质量密度，恰好对应数值 \(1\)。

一般集合上则可能出现尺度振荡。令 \(C\) 为三分 Cantor 集，
\(s=\log2/\log3\)，并记 \(M=\mathcal H^s(C)\in(0,\infty)\)。
相似缩放与测度可加性给出：每个第 \(n\) 层基本区间内的部分都有质量 \(2^{-n}M\)。
在端点 \(0\)，取两列半径
\[
 r_n=3^{-n},\quad R_n=2\cdot3^{-n}.
\]
这两个闭球与 \(C\) 的交集，只相差至多有限个端点：
从最左基本区间的右端到下一基本区间的左端，中间仍是空隙。
单点的 \(\mathcal H^s\) 测度为零，故
\[
 \frac{\mathcal H^s\bigl(C\cap\overline B(0,r_n)\bigr)}{\alpha_s r_n^s}
 =\frac{M}{\alpha_s},
 \quad
 \frac{\mathcal H^s\bigl(C\cap\overline B(0,R_n)\bigr)}{\alpha_s R_n^s}
 =\frac{M}{2^s\alpha_s}.
\]
两个常数不同，所以 \(\Theta^s(C,0)\) 不存在。
这个计算只证明端点 \(0\) 的失败，不能据此声称已经证明
Cantor 集几乎处处的密度都不收敛。

即使是直线段的并，也不能把密度的几乎处处结论改成处处结论。
在平面中取两条垂直线段的并 \(E\)，交点为各自的内点。
对充分小的 \(r\)，交点附近的总长度为 \(4r\)，而 \(\omega_1r=2r\)，
故交点的一维密度为 \(2\)。
交点本身是 \(\mathcal H^1\)-零测集，因此这并不违反上密度定理。

\subsection{可整流性的预告}
\label{b1:subsec:150504}

到此为止，Hausdorff 测度为任意集合提供了长度、面积及非整数尺度的统一测量办法。
但“具有有限的 \(k\) 维测度”与“可以分成可参数化的 \(k\) 维曲面片”是两种不同的要求。
前者只限制覆盖代价，后者还要求局部形状可以用映射描述。
第\ref{b1:ch:19}章将正式讨论可整流性；这里先说明为什么后面还需要映射的微分。

设 \(f:A\subset\mathbb R^k\rightarrow\mathbb R^d\) 是 Lipschitz 映射。
本章已经知道
\[
 \mathcal H^k\bigl(f(A)\bigr)
 \leq\operatorname{Lip}(f)^k\mathcal H^k(A).
\]
这能控制像集的大小，却不能区分哪些方向被压缩，也不能计算映射多次覆盖同一部分的影响。
如果在某点附近 \(f\) 可以用一个线性映射逼近，
那么线性映射的秩与体积伸缩就有机会进入局部测度的计算。
为把这一想法变成定理，必须先证明 Lipschitz 映射几乎处处可微，
再把点态线性近似整合为面积公式。这分别是第\ref{b1:ch:16}章与第\ref{b1:ch:17}章的任务。

因此，密度是连接覆盖测度与局部几何的一项资料，却不是全部资料。
本节的上密度定理没有产生切平面，也没有提供参数化；
平面例子中的精确密度为 \(1\)，依赖我们事先知道集合位于一个平面。
后续论证会逐步回答何时能够反过来从映射或测度的精细性质恢复这种结构，
而不把这里尚未证明的逆向命题当作现成工具。

% 本章小结（不计入编号）
\begin{chaptersummary}
Hausdorff 测度把集合大小的计算分为两步：先在限定尺度下计算覆盖代价，
再让允许的尺度趋于零。改变直径的幂得到维数的临界指数，
而归一化常数保证在整数维欧氏空间中恢复原有的 Lebesgue 测度。
计算维数时，上界通常来自显式覆盖，下界则来自不容质量过分集中的测度。
Frostman 引理说明，对紧集而言，只要相应 Hausdorff 测度为正，这种测度就能够构造。

密度把这些全局计算带回到点附近。有限 Hausdorff 测度集的上密度几乎处处有正的统一下界和有限上界，
但密度极限及局部平面结构尚需更多条件。下一章将从 Lipschitz 映射的可微性入手，
把距离控制推进到线性近似，为精确计算像集的测度作准备。
\end{chaptersummary}


\BookExercises

% 第十五章习题临时稿；合入章聚合文件已有的 BookExercises 之后。
% 八题均自拟，来源是本章工具的变式；不冒称取自尚未读到全文的教材习题。
% 统一采用 H^s=c_s h^s，H^s_infty 同系数；不另建标题或书签。

\begin{exercise}\label{b1:ex:15-bilipschitz-norms}
设 \(E\subset\mathbb R^d\)，映射 \(\Phi:E\to\mathbb R^m\) 满足
\[
 a|x-y|\le|\Phi(x)-\Phi(y)|\le b|x-y|
 \quad(x,y\in E),
\]
其中 \(0<a\le b<\infty\)。
\begin{enumerate}
\item\label{b1:item:15-bilipschitz-measure}
证明对任意 \(s>0\) 和 \(A\subset E\)，
\[
 a^s\mathcal H^s(A)\le\mathcal H^s\bigl(\Phi(A)\bigr)
 \le b^s\mathcal H^s(A),
\]
且 \(\dim_{\mathrm H}\Phi(A)=\dim_{\mathrm H}A\)。
\item\label{b1:item:15-norm-independent}
证明：在有限维实向量空间上，换用另一个范数不改变任何子集的 Hausdorff 维数。
这是否意味着相应的 Hausdorff 测度数值也不变？
\end{enumerate}
\end{exercise}
% 来源：自拟；把 Lipschitz 单向估计用于逆映射，并连接有限维范数等价。
\begin{solution}
左侧假设保证 \(\Phi\) 单射，且 \(\Phi^{-1}:\Phi(E)\to E\) 为 \(a^{-1}\)-Lipschitz 映射。
分别对 \(\Phi\) 和 \(\Phi^{-1}\) 应用本章的 Hausdorff 测度估计，得到
\[
 \mathcal H^s\bigl(\Phi(A)\bigr)\le b^s\mathcal H^s(A),
 \quad
 \mathcal H^s(A)\le a^{-s}\mathcal H^s\bigl(\Phi(A)\bigr).
\]
因此两边测度同时为零，也同时为无穷，临界指数相同。

给定有限维空间的一组坐标，可将它识别为 \(\mathbb R^d\)。若 \(N\) 为其上的一个范数，则由三角不等式存在 \(B<\infty\)，使 \(N(x)\le B|x|\)，并且
\[
 |N(x)-N(y)|\le N(x-y)\le B|x-y|.
\]
所以 \(N\) 在欧氏单位球面上连续，取得一个严格正的最小值 \(A\)。齐次性给出 \(A|x|\le N(x)\)。
两个范数遂相互控制，恒等映射满足第一问的双边距离估计。
但测度数值可以变化：在 \(\mathbb R\) 上把距离从 \(|x-y|\) 改为 \(2|x-y|\)，每个集合的直径都变为原来的两倍，故 \([0,1]\) 的一维 Hausdorff 测度从 \(1\) 变为 \(2\)。
\end{solution}

\begin{exercise}\label{b1:ex:15-dense-countable-content}
令 \(D=\mathbb Q\cap[0,1]\)。
\begin{enumerate}
\item\label{b1:item:15-dense-countable-zero}
证明对每个 \(s>0\)，都有 \(\mathcal H^s(D)=\mathcal H^s_\infty(D)=0\)。
即使要求覆盖由开区间组成，覆盖代价的下确界仍为零。
\item\label{b1:item:15-closure-content-jump}
证明 \(\mathcal H^1_\infty(\overline D)=\mathcal H^1(\overline D)=1\)，并比较 \(D\) 与 \(\overline D\) 的维数。
指出为何“用覆盖集的闭包替代覆盖集”和“先把被覆盖集合取闭包”是两种不同的操作。
\end{enumerate}
\end{exercise}
% 来源：自拟；区分覆盖集的闭包处理与待测集合的闭包，检验尺度与拓扑信息的差别。
\begin{solution}
将 \(D\) 列为 \(\{q_1,q_2,\ldots\}\)。单点集组成任意细的可数覆盖，其 \(s\) 次直径幂和为零。
若只许开区间，给定 \(\delta,\varepsilon>0\)，可取以 \(q_j\) 为中心、长度为 \(\ell_j\) 的开区间，使
\[
 0<\ell_j<\delta,
 \quad c_s\ell_j^s<\varepsilon 2^{-j}.
\]
这些区间覆盖 \(D\)，总代价小于 \(\varepsilon\)，所以结论不变。

\(\overline D=[0,1]\)。任意可数集合覆盖 \([0,1]\) 时，每个有界非空覆盖集都可放进同直径的闭区间；由 Lebesgue 外测度的次可加性，这些直径之和至少为 \(1\)。
无界覆盖集的代价已为无穷，不影响这个下界。
又因 \(c_1=1\)，用 \([0,1]\) 自身作覆盖即得
\(\mathcal H^1_\infty([0,1])=1\)。正文的归一化定理给出 \(\mathcal H^1([0,1])=1\)。
因此 \(\dim_{\mathrm H}D=0\)，而 \(\dim_{\mathrm H}\overline D=1\)。

对一个给定覆盖逐项取闭包，不改变每项直径，也仍覆盖原集合；但可数个闭包之并未必包含原集合的闭包。
本题中单点覆盖逐项取闭包后仍只有 \(D\)，并没有变成 \([0,1]\) 的覆盖。
\end{solution}

\begin{exercise}\label{b1:ex:15-snowflake-metric}
设 \((X,\rho)\) 是度量空间，\(0<\alpha<1\)，令
\(\rho_\alpha(x,y)=\rho(x,y)^\alpha\)。
\begin{enumerate}
\item\label{b1:item:15-snowflake-measure}
证明 \(\rho_\alpha\) 是与 \(\rho\) 产生相同拓扑的度量。记下标表示所用度量，证明对 \(s>0\)，
\[
 \mathcal H^s_{\rho_\alpha,\delta}(E)
 =\frac{c_s}{c_{\alpha s}}
   \mathcal H^{\alpha s}_{\rho,\delta^{1/\alpha}}(E),
 \quad
 \mathcal H^s_{\rho_\alpha}(E)
 =\frac{c_s}{c_{\alpha s}}\mathcal H^{\alpha s}_{\rho}(E).
\]
\item\label{b1:item:15-snowflake-dimension}
推出 \(\dim_{\mathrm H}(E,\rho_\alpha)=\alpha^{-1}\dim_{\mathrm H}(E,\rho)\)。
说明这怎样给出拓扑相同而 Hausdorff 维数不同的两个度量，并解释 \(\alpha>1\) 时结论不能照搬。
\end{enumerate}
\end{exercise}
% 来源：自拟；逐覆盖核对距离幂变换，并保留本书维数相关的归一化常数。
\begin{solution}
对 \(u,v\ge0\)，不等式 \((u+v)^\alpha\le u^\alpha+v^\alpha\) 与 \(\rho\) 的三角不等式给出 \(\rho_\alpha\) 的三角不等式。
例如固定 \(u>0\) 后，可由 \((1+t)^\alpha-t^\alpha\) 在 \(t>0\) 上递减且不超过 \(1\) 验证前一不等式；\(u=0\) 时直接成立。
其余度量公理显然成立。两种开球满足
\[
 B_{\rho_\alpha}(x,r)=B_\rho(x,r^{1/\alpha}),
\]
所以拓扑相同。

对任意集合 \(U\)，有
\(\operatorname{diam}_{\rho_\alpha}U=(\operatorname{diam}_\rho U)^\alpha\)。
因而两边允许的覆盖完全对应，每项代价分别为
\(c_s(\operatorname{diam}_\rho U)^{\alpha s}\) 和
\(c_{\alpha s}(\operatorname{diam}_\rho U)^{\alpha s}\)，取下确界即得有限尺度恒等式。
再令 \(\delta\downarrow0\)，得到测度恒等式与维数公式。

在 \([0,1]\) 上，通常距离给出维数 \(1\)，而距离 \(|x-y|^\alpha\) 给出维数 \(1/\alpha\)。
这并不违反前题：恒等映射虽然是同胚，却不满足双边 Lipschitz 估计。
若 \(\alpha>1\)，则 \([0,1]\) 上的三个点 \(0,1/2,1\) 满足
\[
 |1-0|^\alpha=1>2(1/2)^\alpha,
\]
此时距离幂甚至不再是度量。
\end{solution}

\begin{exercise}\label{b1:ex:15-holder-graphs}
设 \(f:[0,1]^d\to\mathbb R^k\) 满足
\[
 |f(x)-f(y)|\le L|x-y|^\alpha
 \quad(x,y\in[0,1]^d),
\]
其中 \(0<\alpha\le1\)。令 \(G_f=\{(x,f(x)):x\in[0,1]^d\}\)。证明
\[
 d\le\dim_{\mathrm H}G_f
 \le\min\left\{\frac d\alpha,\ d+k(1-\alpha)\right\}.
\]
特别地，证明每个 Lipschitz 函数的图像维数恰为其定义域维数；同时说明两种上界来自不同的覆盖方式。
\end{exercise}
% 来源：自拟；结合 Hölder 距离控制与同尺度网格覆盖，不使用后续面积公式。
\begin{solution}
投影 \((x,z)\mapsto x\) 是 \(1\)-Lipschitz 映射，且将 \(G_f\) 映满 \([0,1]^d\)，所以维数下界为 \(d\)。

由于定义域有界，映射 \(x\mapsto(x,f(x))\) 也是 \(\alpha\) 次 Hölder 连续的。
一般地，若 \(T\) 满足 \(|T(x)-T(y)|\le C|x-y|^\alpha\)，则一个直径为 \(r_j\) 的覆盖集，其像的直径不超过 \(Cr_j^\alpha\)。
于是当 \(q>\dim_{\mathrm H}E/\alpha\) 时，可用 \(\alpha q\) 次直径幂和任意小的细覆盖证明 \(\mathcal H^q(T(E))=0\)。
因此 \(\dim_{\mathrm H}T(E)\le\alpha^{-1}\dim_{\mathrm H}E\)，在本题得到第一个上界。

第二种覆盖同时利用图像的水平与竖直方向。取 \(r=1/N\)，将 \([0,1]^d\) 分成 \(N^d\) 个边长 \(r\) 的小立方体。
在每个小立方体上，\(f\) 的振幅不超过 \(L(\sqrt d\,r)^\alpha\)。
对应的图像部分遂包含于一个水平边长 \(r\)、各竖直边长不超过常数倍 \(r^\alpha\) 的长方体。
当 \(r\le1\) 时，该长方体可由至多 \(C r^{-k(1-\alpha)}\) 个边长 \(r\) 的 \((d+k)\) 维立方体覆盖，其中 \(C\) 与 \(N\) 无关。
故整个图像可用至多 \(C r^{-d-k(1-\alpha)}\) 个直径为 \(\sqrt{d+k}\,r\) 的集合覆盖。
对 \(q>d+k(1-\alpha)\)，这给出
\[
 \mathcal H^q_{2\sqrt{d+k}\,r}(G_f)
 \le Cc_q(d+k)^{q/2}r^{q-d-k(1-\alpha)}\longrightarrow0.
\]
两个上界与投影下界合用即可；\(\alpha=1\) 时上下界均为 \(d\)。
\end{solution}

\begin{exercise}\label{b1:ex:15-zero-dimensional-compact}
令
\[
 K=\left\{\sum_{n=1}^{\infty}\varepsilon_n2^{-n^2}:
       \varepsilon_n\in\{0,1\}\right\}\subset\mathbb R.
\]
\begin{enumerate}
\item\label{b1:item:15-zero-compact-geometry}
证明每个 \(x\in K\) 的上述表示唯一，\(K\) 不可数、紧且无孤立点，但 \(\dim_{\mathrm H}K=0\)。
\item\label{b1:item:15-zero-compact-surjection}
利用唯一表示定义
\[
 F\left(\sum_{n=1}^{\infty}\varepsilon_n2^{-n^2}\right)
 =\sum_{n=1}^{\infty}\varepsilon_n2^{-n}.
\]
证明 \(F\) 是 \(K\) 到 \([0,1]\) 的连续满射，却不满足任何正指数的 Hölder 条件。
\end{enumerate}
\end{exercise}
% 来源：自拟；把固定收缩比 Cantor 构造改为超几何收缩，连接维数、完美性与连续映射。
\begin{solution}
记 \(a_n=2^{-n^2}\)、\(T_n=\sum_{j>n}a_j\)。因为 \(a_{n+1}/a_n\le1/8\)，
\[
 T_n\le\frac87a_{n+1}\le\frac17a_n<a_n.
\]
两列数码若第一次在第 \(n\) 位不同，对应两点之差的绝对值至少为 \(a_n-T_n>0\)，故表示唯一。
二进数码列不可数，故 \(K\) 也不可数。

给定 \(K\) 中任意点列，依次抽选子列，使第一位、第二位等数码逐位固定，再取对角子列。
其前 \(n\) 位与某一数码列最终相同，差值不超过 \(T_n\to0\)，所以此子列收敛于 \(K\) 中的点。
因而 \(K\) 序列紧，从而紧。任一点都可通过改动足够靠后的一位数码得到另一个任意接近的点，所以没有孤立点。

固定前 \(n\) 位后，余下尾项落在长度 \(T_n\) 的区间内，故 \(K\) 可由 \(2^n\) 个这样的区间覆盖。
对任意 \(q>0\)，
\[
 c_q2^nT_n^q
 \le c_q(8/7)^q2^{n-q(n+1)^2}\longrightarrow0.
\]
覆盖尺度也趋于零，所以 \(\mathcal H^q(K)=0\)，得到维数为零。

对每个固定 \(n\)，令
\(\eta_n=\min_{1\le j\le n}(a_j-T_j)>0\)。若 \(x,y\in K\) 且 \(|x-y|<\eta_n\)，两者前 \(n\) 位数码必相同，故 \(|F(x)-F(y)|\le2^{-n}\)。
这证明连续性。\([0,1]\) 中每个数都有二进展开，因而 \(F\) 满射；不同点可能由于二进展开的非唯一性而有相同像，但不影响满射结论。

最后，\(a_n\in K\)，\(F(a_n)=2^{-n}\)，\(F(0)=0\)。对任何 \(\alpha>0\)，
\[
 \frac{|F(a_n)-F(0)|}{|a_n-0|^\alpha}
 =2^{\alpha n^2-n}\longrightarrow\infty,
\]
所以任何正指数的 Hölder 条件都失败。
\end{solution}

\begin{exercise}\label{b1:ex:15-energy-choice}
取 \(0<t<1\)，记 \(\lambda\) 为 \([0,1]\) 上的均匀概率测度。
\begin{enumerate}
\item\label{b1:item:15-atom-energy}
证明若有限正测度 \(\mu\) 满足 \(\mu(\{a\})>0\) 对某个点 \(a\) 成立，则 \(I_t(\mu)=\infty\)。
\item\label{b1:item:15-uniform-energy}
直接计算 \(I_t(\lambda)\)，证明
\[
 I_t(\lambda)=\frac{2}{(1-t)(2-t)}.
\]
\item\label{b1:item:15-nonatomic-infinite-energy}
在 \((0,1)\) 上取密度
\[
 \frac{\mathrm d\nu}{\mathrm dx}(x)=\frac t2x^{t/2-1}.
\]
证明 \(\nu\) 是无原子概率测度，\(\operatorname{supp}\nu=[0,1]\)，但 \(I_t(\nu)=\infty\)。
说明“支集维数大于 \(t\)”与“某个指定测度的 \(t\) 次能量有限”不能混为一谈。
\end{enumerate}
\end{exercise}
% 来源：自拟；比较同一支集上的两种密度，不重复正文有限能量推出维数的证明。
\begin{solution}
若 \(\mu(\{a\})>0\)，则 \((\mu\times\mu)(\{(a,a)\})=\mu(\{a\})^2>0\)，而能量核在 \((a,a)\) 为无穷，故第一问成立。

对均匀测度，Tonelli 定理和对称性给出
\[
 I_t(\lambda)
 =2\int_0^1\int_0^x(x-y)^{-t}\,\mathrm dy\,\mathrm dx
 =\frac2{1-t}\int_0^1x^{1-t}\,\mathrm dx
 =\frac2{(1-t)(2-t)}.
\]
对角线是二维 Lebesgue 零测集，核在该处取无穷值不改变此积分。

所给密度非负，且从 \(0\) 到 \(r\) 的积分为 \(r^{t/2}\)，所以总质量为 \(1\)。
它关于 Lebesgue 测度绝对连续，因而无原子；\([0,1]\) 中每点的任意邻域都有正质量，所以支集为 \([0,1]\)。
令 \(A_n=(2^{-n-1},2^{-n}]\)，其中 \(n\ge0\)。则
\[
 \nu(A_n)=2^{-nt/2}(1-2^{-t/2}).
\]
在 \(A_n\times A_n\) 上，\(|x-y|\le2^{-n}\)，所以
\[
 \iint_{A_n\times A_n}|x-y|^{-t}\,\mathrm d\nu(x)\,\mathrm d\nu(y)
 \ge2^{nt}\nu(A_n)^2=(1-2^{-t/2})^2.
\]
这些直积集合两两不交，非负积分求和便得 \(I_t(\nu)=\infty\)。
因此即使支集的维数为 \(1>t\)，测度过度集中在小尺度上仍可使能量发散；均匀测度与 \(\nu\) 的差别在质量配置，而不在支集。
\end{solution}

\begin{exercise}\label{b1:ex:15-separated-similarities}
设 \(S_1,\ldots,S_m:\mathbb R^d\to\mathbb R^d\) 是相似映射，收缩比分别为 \(r_i\in(0,1)\)，其中 \(m\ge2\)。
设非空紧集 \(K\) 满足 \(K=\bigcup_{i=1}^mS_i(K)\)，且这些像两两不交。
令 \(s>0\) 为方程 \(\sum_{i=1}^mr_i^s=1\) 的唯一解。
\begin{enumerate}
\item\label{b1:item:15-similar-natural-measure}
构造 \(K\) 上的概率测度 \(\mu\)，使对每个有限字 \(w=(i_1,\ldots,i_n)\)，
\[
 \mu\bigl(S_{i_1}\circ\cdots\circ S_{i_n}(K)\bigr)
 =(r_{i_1}\cdots r_{i_n})^s.
\]
\item\label{b1:item:15-similar-ball-growth}
证明存在 \(A,B>0\)，使
\[
 Ar^s\le\mu\bigl(B(x,r)\bigr)\le Br^s
 \quad(x\in K,\ 0<r\le\operatorname{diam}K).
\]
进而证明 \(0<\mathcal H^s(K)<\infty\)、\(\dim_{\mathrm H}K=s\)。
\end{enumerate}
\end{exercise}
% 来源：自拟；从单一 Cantor 比例推广到不同相似比，关键新增步骤为停止字与球内有限重数。
% 只假定存在满足不交不变关系的紧集，不另调用迭代函数系统的不动点理论。
\begin{solution}
\(\sum_i r_i^u\) 关于 \(u\ge0\) 连续严格递减，从 \(m\) 降到零，故 \(s\) 存在且唯一。
记 \(p_i=r_i^s\)，先把 \([0,1)\) 按比例 \(p_1,\ldots,p_m\) 分成左闭右开的相邻区间 \(I_1,\ldots,I_m\)。
再把每个 \(I_w\) 按同样比例分成半开区间 \(I_{w1},\ldots,I_{wm}\)，逐级进行。
对有限字 \(w=(i_1,\ldots,i_n)\)，记 \(S_w=S_{i_1}\circ\cdots\circ S_{i_n}\)，所得区间满足
\[
 m_1(I_w)=p_{i_1}\cdots p_{i_n}=:p_w.
\]
半开约定保证每个 \(t\in[0,1)\) 在每一级恰落入一个区间，因而确定唯一的无限数码列。
记其长度为 \(n\) 的前缀为 \(w_n(t)\)。固定 \(x_*\in K\)，令
\[
 \Phi_n(t)=S_{w_n(t)}(x_*).
\]
每个 \(\Phi_n\) 在有限个半开区间上分别取常值，故 Borel 可测。
紧集 \(S_{w_n(t)}(K)\) 递减，其直径至多为 \(r_{\max}^nD\)，其中 \(r_{\max}=\max_i r_i<1\)、\(D=\operatorname{diam}K>0\)。
它们的交恰含一点 \(\Phi(t)\)，且 \(|\Phi_n(t)-\Phi(t)|\le r_{\max}^nD\)。
因此 \(\Phi\) 是这些 Borel 映射的一致极限，也 Borel 可测，且取值于 \(K\)。

定义 \(\mu=\Phi_\#(m_1|_{[0,1)})\)。这是集中在 \(K\) 上的 Borel 概率测度。
对任一有限字 \(w\)，若 \(t\in I_w\)，则 \(\Phi(t)\in S_w(K)\)。
若 \(t\notin I_w\)，其同长度前缀与 \(w\) 不同；由第一层分块不交及相似映射的单射性，同一级不同分块也不交，所以 \(\Phi(t)\notin S_w(K)\)。
于是 \(\Phi^{-1}\bigl(S_w(K)\bigr)=I_w\)，从而
\[
 \mu\bigl(S_w(K)\bigr)=m_1(I_w)=p_w=(r_{i_1}\cdots r_{i_n})^s.
\]
第一问的测度由此直接构成，无须使用可数乘积测度。

记 \(r_*=\min_i r_i\)、\(\Delta=\min_{i\ne j}\operatorname{dist}(S_i(K),S_j(K))>0\)。
对一个字 \(w\)，记 \(S_w\) 的收缩比为 \(r_w\)，去掉末位的字为 \(w^-\)，空字收缩比取 \(1\)。
给定 \(0<r<D\)，沿每条分支在第一次满足 \(r_wD\le r\) 时停止。
所得停止字两两没有前缀关系，覆盖 \(K\)，并满足
\[
 \frac{r_*r}{D}<r_w\le\frac rD.
\]
若 \(w,v\) 是两个不同停止字，令 \(u\) 为它们的最长公共前缀。两块分别落在 \(S_u\) 下不同的第一层分块内，所以
\[
 \operatorname{dist}\bigl(S_w(K),S_v(K)\bigr)
 \ge r_u\Delta>\frac{r_*\Delta}{D}r.
\]
因此，在每个与 \(B(x,r)\) 相交的停止块内任选一个交点，所得点落在 \(B(x,r)\) 内且两两相距大于 \(cr\)，其中 \(c=r_*\Delta/D>0\)。
以这些点为中心、半径 \(cr/2\) 的球两两不交，且均含于 \(B(x,(1+c/2)r)\)。
比较 Lebesgue 体积，知相交停止块数不超过一个只依赖 \(c,d\) 的常数 \(N\)。
每块质量至多为 \((r/D)^s\)，故 \(\mu\bigl(B(x,r)\bigr)\le N(r/D)^s\)。在 \(r=D\) 时用总质量 \(1\) 即可延续这个上界。

为证下界，对给定 \(x\in K\) 沿其唯一编码在直径首次不超过 \(r/2\) 时停止。
所得块包含 \(x\)，并全部落在开球 \(B(x,r)\) 中；其质量大于 \((r_*r/(2D))^s\)。
于是取得双边球估计。若球中心不在 \(K\) 上但球与 \(K\) 相交，取一个交点为新中心，将半径扩大到 \(2r\)，即可得到所有足够小球的同型上界。
由质量分布原理，\(\mathcal H^s(K)>0\)。

另一方面，所有长度为 \(n\) 的字给出尺度趋于零的有限覆盖，且
\[
 c_s\sum_{|w|=n}\bigl(\operatorname{diam}S_w(K)\bigr)^s
 =c_sD^s\left(\sum_{i=1}^mr_i^s\right)^n=c_sD^s.
\]
所以 \(\mathcal H^s(K)\le c_sD^s<\infty\)，临界维数即为 \(s\)。
\end{solution}

\begin{exercise}\label{b1:ex:15-cantor-product}
令 \(C\subset[0,1]\) 为三分 Cantor 集，\(s=\log2/\log3\)，\(k\ge1\) 为整数。
证明
\[
 0<\mathcal H^{s+k}\bigl(C\times[0,1]^k\bigr)<\infty,
 \quad
 \dim_{\mathrm H}\bigl(C\times[0,1]^k\bigr)=s+k.
\]
证明中分别写清楚：下界所用乘积测度的球估计，以及上界所用同尺度覆盖的数量。
不能仅凭“直积维数等于维数之和”作为依据。
\end{exercise}
% 来源：自拟；用规则 Cantor 覆盖和乘积概率建立一个可完全计算的非整数维直积。
\begin{solution}
取 \(C\) 上自然概率测度 \(\mu\)，使每个第 \(n\) 级基本区间的质量为 \(2^{-n}\)。
由正文的 Cantor 估计，或由\cref{b1:ex:15-separated-similarities}应用于 \(x\mapsto x/3\)、\(x\mapsto(x+2)/3\)，存在 \(A<\infty\)，使所有中心及足够小的半径满足
\[
 \mu\bigl((x-r,x+r)\bigr)\le Ar^s.
\]
令 \(\lambda_k\) 为 \([0,1]^k\) 上的 Lebesgue 概率测度，并取乘积概率 \(\nu=\mu\times\lambda_k\)。
\(\mathbb R^{k+1}\) 中半径为 \(r\) 的球包含于它的各坐标投影的直积，而后 \(k\) 个坐标的投影包含在一个边长 \(2r\) 的立方体内，故
\[
 \nu\bigl(B((x,y),r)\bigr)
 \le\mu\bigl((x-r,x+r)\bigr)(2r)^k
 \le A2^kr^{s+k}.
\]
由质量分布原理得到 \(\mathcal H^{s+k}(C\times[0,1]^k)>0\)。

再看上界。用 \(2^n\) 个长度为 \(3^{-n}\) 的基本区间覆盖 \(C\)，并将 \([0,1]^k\) 分成 \(3^{nk}\) 个边长为 \(3^{-n}\) 的立方体。
它们的直积给出 \(2^n3^{nk}\) 个直径为 \(\sqrt{k+1}\,3^{-n}\) 的集合覆盖。
因为 \(3^s=2\)，其总代价为
\[
 c_{s+k}2^n3^{nk}
 \bigl(\sqrt{k+1}\,3^{-n}\bigr)^{s+k}
 =c_{s+k}(k+1)^{(s+k)/2},
\]
与 \(n\) 无关。令 \(n\to\infty\)，即得有限上界及维数等式。
这里的上界依赖 Cantor 集在每一级都有可控数量的同尺度覆盖；并没有对任意两个集合使用一般的直积维数上界。
\end{solution}
